%Document Class
\documentclass[a4paper,12pt]{article}


%Packages 
\usepackage{amscd,amssymb,amsmath,amsfonts,amsthm, mathrsfs}
\usepackage{stmaryrd, accents}
\usepackage{graphicx} % For inserting images
\usepackage{subcaption}
\usepackage{verbatim,enumerate,paralist}
\usepackage{textcomp}
\usepackage{tensor}
%\usepackage[authoryear]{natbib}
\usepackage[colorlinks]{hyperref}
\hypersetup{%
  urlcolor=blue,linkcolor=blue,citecolor=blue
}
\usepackage{pdfpages}
%\usepackage{pdfsync}%  remove this when finalised

\usepackage{tikz-cd}
\usepackage[usenames,dvipsnames]{pstricks}

\usepackage[T1]{fontenc}
%\usepackage[all,2cell,poly,knot,tips]{xy}
%\UseAllTwocells
%\CompileMatrices

% For script lower case
\usepackage[scr=rsfs,cal=boondox]{mathalfa}

%\usepackage[scr=esstix,cal=boondox]{mathalfa}
%%%% Xypic

\usepackage[arrow, matrix, tips, curve, graph, rotate]{xy}
\SelectTips{cm}{10}
\newcommand{\cd}[2][]{\vcenter{\hbox{\xymatrix#1{#2}}}}

%%% This is a patch that Ross Moore sent me to fix a bug which breaks
%%% the \sh macro below.

\makeatletter
\def\matrixobject@{%
  \edef \next@{={\DirectionfromtheDirection@ }}%
  \expandafter \toks@ \next@ \plainxy@
  \let\xy@@ix@=\xyq@@toksix@
  \xyFN@ \OBJECT@}
\let\xy@entry@@norm=\entry@@norm
\def\entry@@norm@patched{%
  \let\object@=\matrixobject@
  \xy@entry@@norm }
\AtBeginDocument{\let\entry@@norm\entry@@norm@patched}
\makeatother

\newcommand{\twocong}[2][0.5]{\ar@{}[#2] \save ?(#1)*{\cong}\restore}
\newcommand{\twoeq}[2][0.5]{\ar@{}[#2] \save ?(#1)*{=}\restore}
\newcommand{\ltwocell}[3][0.5]{\ar@{}[#2] \ar@{=>}?(#1)+/r 0.2cm/;?(#1)+/l 0.2cm/^{#3}}
\newcommand{\rtwocell}[3][0.5]{\ar@{}[#2] \ar@{=>}?(#1)+/l 0.2cm/;?(#1)+/r 0.2cm/^{#3}}
\newcommand{\utwocell}[3][0.5]{\ar@{}[#2] \ar@{=>}?(#1)+/d  0.2cm/;?(#1)+/u 0.2cm/_{#3}}
\newcommand{\dtwocell}[3][0.5]{\ar@{}[#2] \ar@{=>}?(#1)+/u  0.2cm/;?(#1)+/d 0.2cm/^{#3}}
\newcommand{\ultwocell}[3][0.5]{\ar@{}[#2] \ar@{=>}?(#1)+/dr  0.2cm/;?(#1)+/ul 0.2cm/^{#3}}
\newcommand{\urtwocell}[3][0.5]{\ar@{}[#2] \ar@{=>}?(#1)+/dl  0.2cm/;?(#1)+/ur 0.2cm/^{#3}}
\newcommand{\dltwocell}[3][0.5]{\ar@{}[#2] \ar@{=>}?(#1)+/ur  0.2cm/;?(#1)+/dl 0.2cm/^{#3}}
\newcommand{\drtwocell}[3][0.5]{\ar@{}[#2] \ar@{=>}?(#1)+/ul  0.2cm/;?(#1)+/dr 0.2cm/^{#3}}

\newcommand{\pushoutcorner}[1][dr]{\save*!/#1+1.2pc/#1:(1,-1)@^{|-}\restore}
\newcommand{\pullbackcorner}[1][dr]{\save*!/#1-1.2pc/#1:(-1,1)@^{|-}\restore}
\newcommand{\sh}[2]{**{!/#1 -#2/}}

\makeatletter
\newcommand{\xRightarrow}[2][]{\ext@arrow 0359\Rightarrowfill@{#1}{#2}}
\makeatother

\def\blue{\color{blue}}
\def\red{\color{red}}
\def\black{\color{black}}

%\usepackage[T1]{fontenc}
%\usepackage[all,poly,knot,tips]{xy}
%%\usepackage{xypdf} % may not be necessary
%\setlength{\marginparwidth}{2cm}
%\newcommand{\missing}[1]{\marginpar{\color{red}missing:\\#1}}


%Theoretic Environment Setup
\newtheorem{theorem}{Theorem}[section]
\newtheorem{corollary}[theorem]{Corollary}
\newtheorem{lemma}[theorem]{Lemma}
\newtheorem{proposition}[theorem]{Proposition}
\newtheorem{remark}[theorem]{Remark}
\newtheorem{definition}[theorem]{Definition}
\newtheorem{example}[theorem]{Example}
\newtheorem{conjecture}[theorem]{Conjecture}



% use \newtheoremstyle to underline the header
\newtheoremstyle{step}{2\bigskipamount}{\medskipamount}{\upshape}{}{\itshape}{. }{ }{\underline{Step~\thestep}}
\theoremstyle{step}
\newtheorem{step}{Step}[subsection]
\renewcommand{\thestep}{\arabic{step}}

\newcommand{\lra}{\longrightarrow}
\newcommand{\lla}{\longleftarrow}
\newcommand{\Lra}{\Longrightarrow}
\newcommand{\Ra}{\Rightarrow}
\newcommand{\La}{\Leftarrow}
\newcommand{\Lla}{\Longleftarrow}
\newcommand{\ldual}[1]{\mathord{{\let\nolimits\relax\sideset{^\wedge}{}{#1}}}}
\newcommand{\laction}[2]{\mathord{{\let\nolimits\relax\sideset{^{#1}}{}{#2}}}}
\newcommand{\conj}[2]{\mathord{{\let\nolimits\relax\sideset{^{#1}}{}{#2}}}}
\newcommand{\xylongto}[2][2pt]{{\UseTips{}\xymatrix @C+#1{ \ar[r]^{#2}&}}}
\newcommand{\xra}{\xrightarrow}
\newcommand{\xla}{\xleftarrow}
\makeatletter
\newcommand{\xRa}[2][]{\ext@arrow 0359\Rightarrowfill@{#1}{#2}}
\makeatother

% Yoneda
\usepackage[utf8]{inputenc}
\DeclareFontFamily{U}{min}{}
\DeclareFontShape{U}{min}{m}{n}{<-> udmj30}{}
\newcommand\yo{\!\text{\usefont{U}{min}{m}{n}\symbol{'207}}\!}

%% Script Shortcuts 
\def\CA{{\mathscr A}}
\def\CB{{\mathscr B}}
\def\CC{{\mathscr C}}
\def\CD{{\mathscr D}}
\def\CE{{\mathscr E}}
\def\CF{{\mathscr F}}
\def\CG{{\mathscr G}}
\def\CH{{\mathscr H}}
\def\CI{{\mathscr I}}
\def\CJ{{\mathscr J}}
\def\CK{{\mathscr K}}
\def\CL{{\mathscr L}}
\def\CM{{\mathscr M}}
\def\CN{{\mathscr N}}
\def\CO{{\mathscr O}}
\def\CP{{\mathscr P}}
\def\CQ{{\mathscr Q}}
\def\CR{{\mathscr R}}
\def\CS{{\mathscr S}}
\def\CT{{\mathscr T}}
\def\CU{{\mathscr U}}
\def\CV{{\mathscr V}}
\def\CW{{\mathscr W}}
\def\CX{{\mathscr X}}
\def\CY{{\mathscr Y}}
\def\CZ{{\mathscr Z}}
\def\dd{{\colon}}
\def\ot{{\otimes}}

% import these fonts by hand here to avoid clashes with usual blackboard bold usage.
%\pdfmapfile{+bbold.map}
\newcommand{\bbefamily}{\fontencoding{U}\fontfamily{bbold}\selectfont}
\newcommand{\textbbe}[1]{{\bbefamily #1}}
\DeclareMathAlphabet{\mathbbe}{U}{bbold}{m}{n}

% Macro to typeset part of a math formula at a bigger size
\def\makebigger#1#2#3{\hbox{#1$\mathsurround=0pt #2{#3}$}}
\def\bigger#1#2{{\relax\mathpalette{\makebigger#1}{#2}}}

% for example you could use this to define the following macro:
\newcommand{\otb}{{\bigger\large{\mathbbe{\otimes}}}}

\def\theequation{{\thesection.\arabic{equation}}}


\begin{document}

\author{Ross Street}

\title{The natural transformation in mathematics}
\date{\today}
\maketitle
%\noindent {\small{\emph{2020 Mathematics Subject Classification:} 18M10; 18N10; 18G35; 18G90; 18M90}}
%\\
%{\small{\emph{Key words and phrases:} homology; cocompletion; left exact; iso-inserter; subequalizer; coinverter}}

%\tableofcontents
%\newpage
%\begin{abstract}
%\noindent 
\section*{Purpose}

The terms category, functor and natural transformation were coined in the 1945 paper 
\cite{EilMac}. The study that followed was called Categorical Algebra by the time of the
LaJolla conference in 1965 which was attended by G.M. (Max) Kelly from the University of Sydney;
see \cite{EilKel1966}. The subject soon became ``Category Theory'' and Australia has
played a significant role ever since. By now it is common for non-mathematicians to have
heard of the subject since it is often mentioned in connection with its applications. 
 
As requested, in this paper I shall discuss categories, enriched and higher dimensional aspects, and some connections with computer science and theoretical physics. I must admit
that my attempt to do this is a bit technical and longer than intended but I hope the account can be read by anyone who only wants a rough understanding and so can breeze over the precise detail. 

\section{Objective versus Abstract}\label{OvA}

The famous mathematician Hassler Whitney (1907--1989) was awarded the National Medal of Science by the President of the USA in 1980. He spent his final years studying how young children learn mathematics and why older students have mathematics anxiety. When he visited Macquarie University, he was not there to lecture on differential geometry, rather, to see how our Numeracy Centre was helping students overcome their fear of mathematics. 

Over coffee we discussed how children already understood {\em bijections}, that is, one-to-one correspondences, before they could count. He had watched his granddaughter set the dining table: ``A fork for Mommy, a fork for Daddy, one for Grandpa, one for Grandma, one for my brother, and one for me''. She had set up a bijection between diners and forks.  
Number $6$ was not mentioned! 

The point is that, not only are sets (for example, the set of diners and the set of forks in question) primitive concepts, so are bijections between sets. At a stretch, we could say the child was constructing an {\em injective function} $f : D\to F$ from the set $D$ of diners to the set $F$ 
of forks in the draw, injective since no two diners get the same fork. Diners are the {\em elements}
of $D$ and forks are the elements of $F$. 

The arrow notation $f : X \to A$ (or $X\xra{f} A$) for functions was popularised by Witold Hurewicz around 1940.
Here $X$ and $A$ are sets and $f$ is the name of a rule providing an element $f(x)$ of $A$
(the {\em value} of $f$ at $x$) to every element $x$ of $X$.  
It is very important for category theory that to specify the function 
we must provide the rule, the {\em domain} (or {\em source}) set $X$, and the {\em codomain} (or {\em target}) set $A$. The codomain is allowed to contain more elements than the values; if it does not then the function is called {\em surjective}. Below we will see that good geometric notation (the arrow is an example) is also very important. 

Traditionally people spoke of the ``function'' $\mathrm{cos}(\theta)$ when they meant the rule ``adjacent over hypotenuse in a triangle with both a right angle and an angle of $\theta$ radians''.
We now insist which angles $\theta$ are allowed as the domain and which codomain containing the values is specified.   

For sets with a finite number of elements there is an associated number $\# X$ called its {\em cardinality}; it is the number of elements in the set. The set $\mathbb{N}$ whose elements are natural numbers $0, 1, 2, 3, 4, \dots$ is infinite. We say set $X$ has {\em cardinality less than or equal to} that
of set $Y$ when there is some injective function $f : X\to Y$; we write $\# X\le \# Y$. They have the {\em same cardinality} when there is a bijective function between them; this is the same as saying
$\# X\le \# Y$ and $\# Y\le \# X$ (as proved by Schr\"oder-Bernstein). The set $\mathbb{N}$ and the set $\mathbb{R}$ whose elements are real numbers do not have the same cardinality: $\# \mathbb{N}< \# \mathbb{R}$. 

If we have a bijective function $X\xra{\sigma} Y$ between sets $X$ and $Y$, we can {\em transport}
properties of $X$ to $Y$ by applying $\sigma$. For example, there are many different constructions of the set
$\mathbb{R}$ of real numbers starting with $\mathbb{N}$: as decimals, as Dedekind cuts, as equivalence classes of Cauchy sequences, and so on. Arithmetic operations can be defined for each. There are canonical bijective functions between each pair of these sets and the arithmetic operations transport across them. We begin to see that the elements of a set are not as important as how the set relates, via functions, to other sets. Indeed, the elements $x$ of any set $X$ are in bijection with functions $x : \mathbf{1}\to X$ from the set $\mathbf{1}$ with one element $0$.  

My argument is that numbers are more abstract than sets. When we deal with sets rather than numbers, Lawvere-Schanuel call that the {\em objective} point of view. 

Multiplying sets, which is about labelling rectangles, is also less abstract than multiplying numbers. For sets $X$ and $Y$, their {\em (Cartesian) product} $X\times Y$ is the set whose elements
are pairs $(x,y)$ where $x$ is an element of $X$ and $y$ is an element of $Y$. If $X$ has element $1$, $2$ and $3$, and $Y$ has elements $a$ and $b$, then $X\times Y$ has elements as in the rectangle
\begin{eqnarray*}
\begin{matrix} (1,b) & (2,b) & (3,b) \\
(1,a) & (2,a) & (3,a)
\end {matrix}
\end{eqnarray*}
Quite generally we have the abstraction $\# X\times \# Y = \# (X\times Y)$ which, in the example, amounts to $3\times 2 = 6$. 

In other words, if I say $z$ is an element of $X\times Y$ then it means $z$ has two {\em components}, 
$z_1$ in $X$ and $z_2$ in $Y$. Indeed, not only are the elements of $X\times Y$ pairs but, 
for any set $K$, the functions $K\xra{h} X\times Y$ are in a natural bijection with pairs of 
functions $(h_1,h_2)$ where $K\xra{h_1} X$ and $K\xra{h_2} Y$.  

Addition of sets is a little more complicated but still rather concrete. If the sets $X$ and $Y$ are
disjoint (that is, have no common elements) then we could take $X+Y$ to be the union of the two sets. In case they do have common elements, we need labels, call them $1$ and $2$, 
so that the elements of
$X+Y$ become those of the form $(1,x)$ for all $x$ in $X$ and those of the form $(2,y)$ for all $y$ in $Y$. That is, we take copies of $X$ and $Y$ to make them disjoint and then take the union. 

Now something amazing appears which connects multiplication and addition of sets of all sizes. 
For any set $K$, there is a natural bijection between functions $X+Y\xra{k} K$ and pairs 
$(k_1,k_2)$ of functions where $X\xra{k_1}K$ and $Y\xra{k_2} K$. When $X$ and $Y$ are finite
sets of cardinality $m$ and $n$, respectively, the set of functions from $X$ to $Y$ has cardinality
$n^m$. For this reason, we write $Y^X$ for the set of functions from any $X$ to any $Y$. So what we
have is a natural bijection $\theta : K^{X+Y}\to K^X\times K^Y$ defined by $\theta(k) = (k_1,k_2)$.
The corresponding abstract and more limited formula is the familiar numerical identity $\ell^{m+n} = \ell^m\times \ell^n$; this only tells that there is a bijection and not how to construct it. Note also
that the numerical identity $(m\times n)^{\ell} = m^{\ell}\times n^{\ell}$ is an abstraction of the 
natural bijection between $(X\times Y)^K$ and $X^K\times Y^K$.

\begin{example}\label{infnatnumbers} If $X$ is a finite set then there is no bijective function $\mathbf{1}+X\to X$; the cardinalities differ by $1$. However, if $\mathbf{1}\xra{0}\mathbb{N}$ is the first element $0$ of 
$\mathbb{N}$ and $\mathbb{N}\xra{s}\mathbb{N}$ is the {\em successor function} defined by 
$s(n) = n+1$ then the function $\mathbf{1}+\mathbb{N}\xra{t}\mathbb{N}$ with $t_1=0, t_2= s$
is a bijection. (The reader possibly has heard of the {\em Hilbert hotel} with infinitely many rooms.)  
\end{example}

I have used the word ``natural'' several times above in connection with general functions. Category theory began with a paper \cite{EilMac} motivated by the pursuit of a precise definition of this word. Eilenberg and Mac Lane, the authors, formalised the calculus of arrows, yielding the definition of 
``category''. They produced certain arrows between categories which they called ``functors''.
Between functors with the same domain and codomain, they defined {\em higher arrows} called ``natural transformations''.   

Another word that will be used below is ``field''. In mathematics, it means a set with 
addition and multiplication satisfying the usual rules of arithmetic such as the rational numbers 
$\mathbb{Q}$, the real numbers $\mathbb{R}$, the complex numbers of rational numbers 
$\mathbb{C}$, or any of the finite fields such as $\mathbb{F}_7$ the integers modulo $7$. 
However, a ``vector field''
is a space of vectors indicated by arrows of differing length such as in the diagrams in weather reports
showing wind direction and strengths at points on a map. It is the latter which is relevant to physics where
examples are the electric, magnetic and gravitational fields. 

\section{Categories}

It may seem naive to look at sets and functions as forming a very big {\em directed graph} since
so much structure is being ignored. Please bear with me.

A graph $\CG$ has {\em vertices} and {\em edges} joining
the vertices. The graph is directed when each edge is given a direction from one vertex 
(the source) to the other it joins (the target); this direction is indicated by denoting the edges by arrows. 
There is no restriction on how many edges there are between any two vertices and edges can have
the same source as target. This means that the vertices should be labelled and so should the edges.
We will write  $\mathrm{v}\CG$ for the set of vertices. The set of edges from vertex $A$ to vertex $B$ is denoted by $\CG(A,B)$.
A simple example is shown in Figure~\eqref{exgraph} with vertices $A$, $B$, $C$ and directed edges
$e$, $f$, $g$, $h$, $k$.
\begin{eqnarray}\label{exgraph}
\begin{aligned}
\xymatrix{
A \ar@(ul,dl) [] _ {e} \ar@/^/[rr] ^ f
&& B \ar@/^/[ll] ^ {g}
&& C \ar@/^/[ll] ^ {h} \ar@/_/[ll] _ {k}
}
\end{aligned}
\end{eqnarray}
Graphs with infinitely many vertices or edges obviously cannot be drawn in full on a piece of paper,
however, finite subgraphs of interest can be.

There is an important operation on functions that we did not mention in Section~\ref{OvA}.
If they match up appropriately, functions can be substituted one into the other. 
This is called {\em composition}. Explicitly, if $X\xra{f}Y$ and $Y\xra{g}Z$ are functions between
sets then there is a function $X\xra{g\circ f}Z$ defined by $(g\circ f)(x) = g(f(x))$.
This composition is {\em associative} which means, for $X\xra{f}Y$, $Y\xra{g}Z$, $Z\xra{h}W$,
a three-fold composite can be formed two at a time in the two possible ways with the same result;
that is, 
\begin{eqnarray}\label{assoc}
(h\circ g)\circ f = h\circ (g\circ f) \ .
\end{eqnarray} 

Each set $X$ has an {\em identity function} $X\xra{1_X}X$ defined by $1_X(x) = x$. 
It has the unsurprising property that it composes without changing the function you are 
composing with: that is, for any function $X\xra{f}Y$, 
\begin{eqnarray}\label{ident}
1_Y\circ f = f \ , \ f\circ 1_X = f  \ .
\end{eqnarray} 
The sets $\CA(A,B)$ of morphisms are called the {\em hom-sets} of $\CA$ for the historical 
reason that many categories originally of interest were from algebra where the morphisms 
are called ``homomorphisms''. 

\begin{definition}\label{categorydef} A {\em category} is a directed graph $\CA$ in which the vertices are called {\em objects}
and the directed edges are called {\em morphisms}, equipped with functions
\begin{eqnarray}\label{compo}
\circ : \CA(B,C) \times \CA(A,B) \lra \CA(A,C)
\end{eqnarray}
written between the arguments, called {\em composition} for all objects $A, B, C$, and morphisms $1_A : A\to A$ called {\em identities}
for all objects $A$, subject to equations \eqref{assoc} and \eqref{ident} as axioms.    
\end{definition}

\begin{example}
Section~\ref{OvA} is all about the motivating example: the category $\mathrm{Set}$ of sets.
The objects are sets (where some limit on cardinality may be imposed to avoid paradoxes), the morphisms are functions, and the composition and identities are as already explained.
\end{example}

\begin{example}\label{Matex}
Let $\mathbb{F}$ denote a field in the mathematical sense (that is, a set with addition and multiplication
satisfying the usual rules of arithmetic). There is a category $\mathrm{Mat}_\mathbb{F}$ of {\em matrices}.
The objects are natural numbers. The morphisms $a : n\to m$ are $m\times n$ matrices (that is,
rectangular arrays of elements $a_{ij}$ of $\mathbb{F}$ with $m$ rows and $n$ columns). The composite $a\circ b$ of $n\xra{b} m\xra{a} \ell$ is given by matrix multiplication:
\begin{eqnarray*}
(a\circ b)_{ik} = \sum_{j=1}^{m} a_{ij}b_{jk} = a_{i1}b_{1k} + a_{i2}b_{2k}+ \dots + a_{im}b_{mk} \ .
\end{eqnarray*} 
 The identity matrix of $n$ is the {\em Kronecker delta} $\delta$ defined by
\begin{equation*}
\delta_{ij} =
\begin{cases}
1 & \text{if } \ i=j  \\
0 & \text{if } \ i\neq j \ .
\end{cases}
\end{equation*}
\end{example}

\begin{example}\label{freecat}
For every directed graph $\CG$ there is a category $\mathrm{F}\CG$ called the {\em free category on} $\CG$. The objects are the vertices of $\CG$. The morphisms of $\mathrm{F}\CG$ from $A$ to $B$ are {\em (directed) paths} in $\mathrm{F}\CG$; diagram \eqref{pathn} is a path of length $n$.  
\begin{equation}\label{pathn}
A= A_0\xra{f_1} A_1 \xra{f_2} A_2\xra{f_3} \dots \xra{f_n} A_n= B
\end{equation}
We allow the {\em empty path} of no edges from $A$ to $A$; its length is $0$.
Composition of paths is by concatination so that lengths are added. The empty paths are the identity morphisms.
\end{example}

Let me pause the examples at this point to mention something which is the low level example of a much more complex aspect for higher categories. Suppose the path \eqref{pathn} consists of morphisms in
a category $\CA$, not just edges in any graph. We can compose the morphisms two at a time in lots of different ways to obtain a morphism from $A$ to $B$; for example, we could start with $f_2\circ f_1$
and continue to the right along the path to give 
\begin{eqnarray}\label{rtbrack}
f_n\circ (f_{n-1}\circ \dots \circ (f_3\circ (f_2\circ f_1))\dots ) : A\to B 
\end{eqnarray}
in $\CA$. It is possible to prove a {\em Generalised Associativity Law} which states that the two axioms
\eqref{assoc} and \eqref{ident} imply that all ways of inserting meaningful brackets in the path result in
the same iterated composite $A\to B$. A proof relevant to symbolic computation is to replace the three
equalities of \eqref{assoc} and \eqref{ident} by rewriting rules which allow occurrences of the left hand sides to be rewritten as the right hand sides, in any path, to obtain a new path. With this, brackets are moving further to the right and the number of identities is being reduced. This stops when all the bracketing is from the right as in \eqref{rtbrack} with no identities unless $A=B$ and the result is $1_A$.
This termination where no left hand sides occur anymore is called the {\em normal form} for the rewriting system, and results in the unique composite for the path in $\CA$. 

A diagram in a category that involves various paths of morphisms is said to {\em commute} when
any two paths between the same objects in the diagram have the same composite. For example,
the diagram \eqref{commdiagex} commutes if and only if the four equations \eqref{4equns} hold. 
\begin{eqnarray}\label{commdiagex}
\begin{aligned}
\footnotesize{\xymatrix{
C \ar[dd]_-{e} \ar[rd]_-{f}  \ar[rr]_-{g} & & D \ar[rr]_-{h} \ar[rd]_-{i} & & E \ar[dd]_-{j}  \\
& F \ar[ld]^-{k} \ar[rr]_-{l} & & G  \ar[ldld]^-{m} \\
H \ar[rd]_{n} & & & & I \ar[d]_-{p} \\
& K  \ar[rrr]^-{q} & & & L } }
\end{aligned}
\end{eqnarray}
\begin{eqnarray}\label{4equns}
i\circ g = \ell\circ f \ , \ p\circ j\circ h = q \circ m \circ i \ , \ k \circ f = e \ , \ m \circ \ell = n\circ k 
\end{eqnarray}




\begin{example} \label{discretecat} Every set $X$ can be regarded as the set of vertices of a graph with no edges; then $\mathrm{F}X$ is called
the {\em discrete category on $X$} where the only morphisms are identities.
\end{example}
\begin{example} \label{MonoidDef}
If a category $\CA$ has only one object $\star$, the category can be identified with a {\em monoid}; that is, a set $M$ with an associative binary operation~$\cdot$ having an identity $1$. The connection is $\CA(\star,\star) = M$. We already saw this in Example~\ref{onevertex} where $\CA=\mathrm{F}\CG$ and $M= \mathbb{N}$ (the binary operation is addition).
A slogan is that ``a category is just a monoid spread over several objects'', but ``several'' can mean lots.  
\end{example}
\begin{example}\label{onevertex} 
If a graph $\CG$ has only one vertex, say $\star$, and one edge $\star\xra{e} \star$ then the morphisms of $\mathrm{F}\CG$ are paths $\star\xra{e} \star \xra{e} \star\xra{e} \dots \xra{e} \star$ which are
totally known when we know the length. So we have the identification 
$$\mathrm{F}\CG(\star,\star) = \mathbb{N} \ ;$$ 
Composition becomes addition of natural numbers.
\end{example}
\begin{example}\label{alphabet} More generally, if the graph $\CG$ has only one vertex $\star$ then we can think of the edges $\star\xra{a} \star$
as letters in some {\em alphabet} $\Lambda$. The morphisms in $\mathrm{F}\CG$ are then called {\em words} in the alphabet $\Lambda$. The monoid of words is denoted by $\Lambda^*$.    
\end{example}

\begin{example} There is a category $\mathrm{Mon}$ whose objects are monoids. The morphisms
$f : M\to N$ are functions which preserve the binary operation and identity: $f(x\cdot y) = f(x)\cdot f(y)$,
$f(1) = 1$. Composition and identities are as for functions: if you compose two monoid morphisms it is again a monoid morphism. 
\end{example}

\begin{example} Every directed graph $\CG$ gives an {\em opposite graph} $\CG^{\mathrm{op}}$
obtained by keeping the same objects and reversing all the directions of the edges. This applies to
categories $\CA$ to give an {\em opposite category} $\CA^{\mathrm{op}}$ where the composition is obtained by reversing the one in $\CA$. While this might seem a very formal process it can happen
that if $\CA$ is a known category coming from algebra then $\CA^{\mathrm{op}}$ can be identified
with another known category, usually more geometrical! An example is that $\mathrm{Set}^{\mathrm{op}}$ can be identified with the category of complete atomic Boolean algebras.   
\end{example}

\section{Product and Exponential}\label{ProdExp}
The category $\mathrm{Set}$ of sets can be used to suggest many notions for a general category $\CA$ if the notion can be expressed in terms of the data at hand: objects, morphisms and composition.

A {\em left inverse} for a morphism $S\xra{m}A$ is a morphism $A\xra{\ell}S$ such that $\ell\circ m = 1_S$. In $\mathrm{Set}$, a function with non-empty domain has a left inverse if and only if it is injective. A {\em right inverse} for a morphism $A\xra{e}C$ is a morphism $C\xra{r}A$ such that $e\circ r = 1_C$. In $\mathrm{Set}$, a function has a right inverse if and only if it is surjective; this is equivalent to saying our set theory satisfies the {\em Axiom of Choice}. If a morphism $A\xra{f}B$
has both a left inverse $\ell$ and a right inverse $r$ then
$$\ell = \ell \circ 1_B = \ell \circ (f\circ r) = (\ell \circ f)\circ r = 1_A\circ r = r \ .$$    
That is, the left and right inverse for $f$ are equal if they exist; then we say $f$ is {\em invertible}. Notice how the proof uses the 
associativity and identity properties of composition in $\CA$. This means we can talk about {\em the}
inverse and denote it by $f^{-1}$. The invertible morphisms in $\mathrm{Set}$ are the bijective functions. Two objects $A, B$ in $\CA$ behave in the same way when there is an invertible function between them; we say the objects are {\em isomorphic} and we write $A\cong B$.
Objects of $\mathrm{Set}$ are isomorphic when they have the same cardinality. 

\begin{example} A category in which all morphisms are invertible is called a {\em groupoid}.
If a groupoid only has one object, and so is a monoid (see Example~\ref{MonoidDef}), the monoid is called a {\em group}.
\end{example} 

An object $T$ of a category $\CA$ is {\em terminal} (or {\em a terminator}) when, for all objects
$A$ of $\CA$, there is precisely one morphism $A\to T$. If we have that there exists such
a morphism $A\to T$ but not necessarily unique, the object is said to be {\em weakly terminal}. Any two terminators are isomorphic. Some categories do not have terminators at all.
 The terminal objects of $\mathrm{Set}$
are those with precisely one element, such as the set we have denoted by $\mathbf{1}$. 
Sometimes we also denote a chosen terminator in other categories by $1$. In
$\mathrm{Set}$, the morphisms $\mathbf{1}\to A$ are the elements of the set $A$ and so
determine $A$ up to isomorphism. In a given category $\CA$, all the morphisms $T\to A$  
from a terminator $T$ do not necessarily determine the object $A$. A terminal object in
$\CA^{\mathrm{op}}$ is called an {\em initial} object. The empty set is initial in $\mathrm{Set}$.

Suppose $\CA$ is a category with a terminator $1$. Consider the category $\mathrm{NN}_{\CA}$
whose objects are triplets $(X, 1\xra{x_0}X, X\xra{u}X)$ of data in $\CA$ and whose morphisms
$(X,x_0,u)\xra{f}(Y, y_0,v)$ are morphisms in $\CA$ such that $f\circ x_0=y_0$ and $f\circ u=v$;
composition is as in $\CA$. An initial object in $\mathrm{NN}_{\CA}$ is called a 
{\em Peano-Lawvere natural numbers object (NNO) in $\CA$} and denoted by $(N, 0, s)$.
In other words, given any $(X,x_0,u)$, there is a unique morphism $r$ such that directed 
paths in the diagram \eqref{(w)nno} with the same domain and codomain are equal when composed
(we say ``the diagram commutes''). 
\begin{eqnarray}\label{(w)nno}
\begin{aligned}
\xymatrix{
1 \ar[r]^-{0} \ar[rd]_-{x_0} & N \ar[d]^-{r} \ar[r]^-{s}  & N  \\
& X \ar[ru]_-{u} }
\end{aligned}
\end{eqnarray}
In $\mathrm{Set}$, any NNO $N$ is isomorphic to the $\mathbb{N}$ of Example~\ref{infnatnumbers};
then the initiality condition is a recursive (inductive) definition of sequences $r$ in $X$ 
when we ask that $r(0) = x_0$ and $r(n+1)= u(r(n))$. 
If the uniqueness condition is deleted in the definition of NNO, we obtain the definition of {\em weak} NNO. 

A {\em product} of two objects $A$ and $B$ in a category $\CA$ is defined as a terminal object
in the category $\mathrm{Spn}_{\CA}(A,B)$ of {\em spans} from $A$ to $B$. A span from $A$
to $B$ is a diagram $A\xla{s_1}S\xra{s_2}B$ in $\CA$. A morphism from span $A\xla{s_1}S\xra{s_2}B$ to span $A\xla{t_1}T\xra{t_2}B$ is a morphism $S\xra{f}T$ in $\CA$ such that $s_1 = t_1 \circ f$
and $s_2 = t_2 \circ f$. Composition is that of $\CA$. It follows that the product is unique up to isomorphism in $\mathrm{Spn}_{\CA}(A,B)$. We denote a particular choice of the product by
\begin{eqnarray*}
A\xla{\mathrm{pr}_1}A \times B\xra{\mathrm{pr}_2}B
\end{eqnarray*}
and refer to $A \times B$ as {\em the} product of $A$ and $B$ with {\em projections} $\mathrm{pr}_1$,
$\mathrm{pr}_2$ in $\CA$. Spans $A\xla{s_1}S\xra{s_2}B$ are in bijection with morphisms 
$S\xra{s} A\times B$ where $\mathrm{pr}_1\circ s = s_1$ and $\mathrm{pr}_2\circ s = s_2$.
In a given category products may not exist; they do in $\mathrm{Set}$ where they agree with the
Cartesian product already discussed. A product in $\CA^{\mathrm{op}}$ is called a {\em coproduct}
or {\em sum} in $\CA$. The coproduct in $\mathrm{Set}$ is disjoint union. 

A category $\CA$ is said to {\em have finite products} when it has a product for every pair of objects
and a terminator. If $U\xra{f}V$ and $A\xra{g}B$ are morphisms in $\CA$ then there is a morphism
$U\times A\xra{f\times g} V\times B$ defined by 
$$\mathrm{pr}_1\circ (f\times g) = f\circ \mathrm{pr}_1
\ \text{ and } \ \mathrm{pr}_2\circ (f\times g) = g\circ \mathrm{pr}_2 \ .$$     

Suppose $\CA$ has finite products. For each pair of objects $Y$, $Z$ in $\CA$, consider the category
$\mathrm{E}_{\CA}(Y,Z)$ whose objects are pairs $(U, U\times Y\xra{h} Z)$ and whose morphisms
from $(U, U\times Y\xra{h} Z)$ to $(V, V\times Y\xra{k} Z)$ are morphisms $U\xra{f}V$ such that 
$h= k\circ (f\times 1_Y)$. Composition is that in $\CA$. A particular choice of terminal object
of $\mathrm{E}_{\CA}(Y,Z)$ is denoted by $(Z^Y, Z^Y\times Y\xra{\mathrm{ev}_Y} Z)$.
We call the object $Z^Y$ the {\em exponential of $Z$ to the power $Y$} and the morphism
{\em evaluation at $Y$}. This means that, for each morphism $U\times Y\xra{h} Z$, there
exists a unique morphism $U\xra{\hat{h}}Z^Y$ such that $h = \mathrm{ev}_Y \circ (\hat{h}\times 1_Y)$.
The process taking $h$ to $\hat{h}$ is often called {\em currying} after the logician H.B. Curry. 
 In $\mathrm{Set}$, the object $Z^Y$ is the set of functions from $Y$ to $Z$ and $\mathrm{ev}_Y$
 is the {\em evaluation} function taking the function $U\times Y\xra{h}Z$ to the function $U\xra{\hat{h}}Z^Y$ defined by $h(u)(y) = h(u,y)$. 
 
 A category $\CC$ is called {\em cartesian closed} when it has finite products and an exponential $Z^Y$ for every pair of objects $Y, Z$. 
 
 \section{Effective calculability}\label{Effcal}
  
The first part of this section is informed by Lambek-Scott's book \cite{LamSco}.

 K. G\"odel, A. Church, S.C. Kleene and A. Turing grappled with the question of which
 numerical functions could be calculated using an effective algorithm. This was the 
 basis of computer programming and started in the 1930s. These functions 
 $\mathbb{N}\xra{f}\mathbb{N}$, called {\em recursive}, are those any computer has any hope
 of calculating by a program. To say what effective algorithms were, Church came up with a 
 kind of formal language called a {\em $\lambda$-calculus} and Turing with a mathematical structure called a {\em deterministic Turing machine}.   
 
 In 1958 at MIT, J. McCarthy created a programming language called LISP based on 
 $\lambda$-calculus which is still widely used in studying {\em artificial intelligence (AI)}. A decade later, J. Lambek \cite{Lam} made the connection between deductive systems (such as the $\lambda$-calculus) and category theory (cartesian closed categories in particular). The connection between {\em typed} $\lambda$-calculus and cartesian closed categories was clarified by D. Prawitz \cite{Pra}, C.R. Mann \cite{Man} and R.A.G. Seely \cite{See}. A typed $\lambda$-calculus $\CL$ consists of {\em types}, {\em terms} of each type, and {\em equations} between terms; equations satisfying a certain closure condition are said to ``hold''. There is a type $N$ of ``natural numbers'' with associated term forming operations for ``successor'' and ``iterator''. In the $\lambda$-calculus, a term can take another function as input; this makes it a {\em higher-order language}. 
   
 There is a category $\mathbf{C}(\CL)$ whose objects are the types of $\CL$. The morphisms
 $A\to B$ are equivalence classes of pairs $(x,\phi(x))$ where $x$ is a variable of type $x$
 and $\phi(x)$ a term of type $B$ whose only free variable is $x$. Free variables are those
 that are not being quantified over by such quantifiers as ``for all $x$'' or ``there exists $x$''.
 The equivalence relation on morphisms involves a slightly technical description of equality. 
 
 It turns out that $\mathbf{C}(\CL)$ is a cartesian closed category containing a weak natural numbers object. In fact, it is the free such in a precise sense that should become clearer below. 
 The higher-order nature of the language $\CL$ is reflected in the exponentials in $\mathbf{C}(\CL)$. We will see that higher-order category theory also comes about from exponentials.   
 
 Of course mathematicians have been considering functions of
 functions as far back as L. Euler (1707 - 1783) who furthermore looked for functions $\phi$ for which 
 $H(\phi)$ was a minimum value of some other function $H$; the subject of {\em calculus of variations}. 
 In 2003, T. Ehrhard and L. Regnier introduced {\em differential $\lambda$-calculus} into functional programming to connect with these areas of mathematical analysis and, by now provides a setting for neuro-symbolic AI. This led to Blute-Cockett-Seely \cite{BluCocSee} creating {\em differential categories} as relevant to functional programming.
More generally, a concept defined by J. Rosick\'y \cite{Ros} in 1984 was revitalised in \cite{BluCocSee2} and called a {\em tangent category}.    

A {\em quantum computer} deals with a single input state and produces a single output state, but in between it can be in multiple superposed states at once. For a theory of quantum computing, there are many variants of a {\em quantum Turing machine}. They involve probabilities which are complex numbers and are not for creating programming languages. The possibility of building real world quantum computers relies on {\em quantum mechanics}. No such computer exists at this point. The problem is to do with background noise
that was not an issue in the classical case. Suggestions for solution involve fault-tolerant quantum computation \cite{LiWil, Kit} using the {\em chain complexes} of algebraic topology.
Chain complexes have a long history in category theory and the monoidal category of such was a major example of a base for enriched category theory (see Section~\ref{enriched}).

Apart from product and coproduct in a category, there can be other kinds of operations on objects of
a category. Because of its relationship with the {\em tensors} used in relativity theory and differential
geometry, one such operation for vector spaces is called the {\em tensor product}. Consequently,
early in the development of category theory, categories equipped with an operation having some of
the properties of tensor product were studied; we follow the terminology of \cite{EilKel1966} and call
them {\em monoidal categories}. Tangent categories are monoidal categories with extra structure.
Also, there are modern approaches to quantum mechanics (see \cite{CoeGog}) based on diagrams developed for monoidal categories. These diagrams are also important in physics and are of a type used earlier by R. Feynman (see \cite{BaeDol2}) and by R. Penrose \cite{Pen1971}.       

\section{Functors}\label{Fun}
Directed graphs form a big category denoted by $\mathrm{Gph}$.
(We will only be dealing with directed graphs so will usually just call them graphs unless emphasis is required.) 
The objects of $\mathrm{Gph}$ are graphs.
A {\em graph morphism} $F : \CG \to \CH$ consists of a function 
$F : \mathrm{v}\CG \to \mathrm{v}\CH$ taking each vertex $A$ of $\CG$ to a vertex $FA$ of 
$\CH$, and, for all pairs $A,B$ of vertices, a function $F_{A,B} : \CG(A,B)\to \CH(FA, FB)$.
Composition and identities are those for functions.
\begin{example}\label{inclgenerators} Recall the free category $\mathrm{F}\CG$ constructed in Example~\ref{freecat}. There is a graph morphism $\CG\to \mathrm{F}\CG$ taking each edge to that edge regarded as a path of length 1. If $\CA$ is a category, there is a functor $\mathrm{F}\CA\to \CA$
taking each path to the unique composite as explained immediately after Example~\ref{freecat}. 
\end{example} 
\begin{example} A graph morphism $F$ from the graph \eqref{graphgraph} to $\mathrm{Set}$
(just regarded as a graph) is precisely a directed graph: $FE$ is the set of edges, $FV$ the set of
vertices, and the functions $Fd$ and $Fc$ assign domain and codomain vertices to edges.  
\begin{eqnarray}\label{graphgraph}
\begin{aligned}
\xymatrix{
V
&& E \ar@/^/[ll] ^ {d} \ar@/_/[ll] _ {c}
}
\end{aligned}
\end{eqnarray}
\end{example}

\begin{definition} A {\em functor} is a graph morphism $T : \CA\to \CB$ which
preserves composition and identities:
$$ T_{A,C}(g\circ f) = T_{B,C}(g)\circ T_{A,B}(f) \ \text{ and } \ T(1_A) = 1_{TA} $$
for all $A\xra{f}B\xra{g}C$ and $A$ in $\CA$.    
\end{definition}
\begin{example} For any object $K$ of a category $\CA$, the functor $\CA\xra{\CA(K,-)} \mathrm{Set}$ taking $A\xra{f}B$ to the function $\CA(K,A)\xra{\CA(K,f)}\CA(K,B)$ whose value
at $K\xra{a}A$ is $K\xra{f\circ a}B$. These are called {\em representable functors}; $K$ is the
representing object.  
\end{example}
\begin{example} Suppose $\CA$ is a category with finite products. For each object $K$ of
$\CA$, there is a functor $\CA\xra{K\times -} \CA$ taking $A\xra{f}B$ to 
$K\times A\xra{1_K\times f}K\times B$.  
\end{example}
\begin{example} Suppose $\CA$ is a cartesian closed category. For each object $K$ of
$\CA$, there is a functor $\CA\xra{(-)^K} \CA$ taking $A\xra{f}B$ to the currying 
$A^K\xra{f^K} B^K$ of $A^K\times K\xra{f\circ \mathrm{ev}_K}B$. There is also a functor
$\CA^{\mathrm{op}}\xra{K^{(-)}} \CA$ taking $A\xra{f}B$ to the currying $K^B\xra{K^f}K^A$
of $K^B\times A \xra{\mathrm{ev}_B\circ (1_{K^B}\times f)}K$.  
\end{example}
\begin{example} For any graph $\CG$ and any category $\CA$, there is a bijection between functors 
$\mathrm{F}\CG\to \CA$ and graph morphisms $\CG\to \CA$ defined by restriction along the graph
morphism of Example~\ref{inclgenerators}.
\end{example}
\begin{example}\label{objmorfunctors} Let us identify the one-element set $\mathbf{1}$ 
with the free category on it. Let $\mathbf{2}$ denote the free category on the graph $0\xra{\mu}1$. Functors $\mathbf{1}\to \CA$ are in bijection with objects of $\CA$. Functors
$\mathbf{2}\to \CA$ are in bijection with morphisms of $\CA$.     
\end{example}
\begin{example}\label{gpreps} A functor $L : G\to \CX$ from a group $G$ (regarded as a one-object category) amounts to an object $X$ together with a function $\ell : G\to \CX(X,X)$ which
takes the group multiplication to composition and takes the identity element of $G$ to $1_X$.
This concept predates the definitions of category and functor. Since functors preserve invertibility, $\ell : G\to \CX(X,X)$ actually lands in the {\em automorphism group} $\mathrm{Aut}_{\CX}X$ of $X$, that is, in the submonoid of $\CX(X,X)$ consisting of the 
invertible morphisms. 
We call $\ell$ an {\em action} of $G$ on $X$ and say $L$ is a {\em representation} of $G$ in $\CX$. Felix Klein (1849-1925) said we could define a ``geometry'' as a representation of a
group in the category of sets; then ``doing geometry'' is studying concepts in $X$ which
are invariant under the action of $G$. For example, studying triangles up to congruence belongs to the geometry where $X$ is the cartesian plane and $G$ is the group
of distance-preserving bijections $X\to X$.     
\end{example}

It is reasonable to call a functor $L : \CA\to \CX$ a {\em representation of $\CA$ in $\CX$}, where $\CA$ is any category, not necessarily a group. Here is a quote from Baez-Dolan \cite{BaeDol1}:

\textit {``One important lesson we have learned from topological quantum field theory is that describing dynamics using group representations is only a special case of describing it using category representations.''} 

Topological quantum field theories themselves are functors with domain categories whose morphisms
are not functions, rather they are geometric objects like curves or surfaces. We will return to this in Sections~\ref{Bklt}~and~\ref{Frobenius}.

\section{Natural transformations}\label{NatTrans}

There is a big category $\mathrm{Cat}$ of categories. The objects are categories.
The morphisms $T : \CA\to \CB$ are functors. Composition is as for graph morphisms.

\begin{theorem} The (big) category $\mathrm{Cat}$ is cartesian closed. 
\end{theorem}
A terminal object is any set with one element such as the set $\mathbf{1}$. 

The product $\CA\times \CB$ of two categories is what you might expect: an object is a pair $(A,B)$ of objects $A$ in $\CA$ and $B$ in $\CB$ while a morphism
$(A,B)\xra{(f,g)}(C,D)$ is a pair of morphisms $f$ in $\CA$ and $g$ in $\CB$. 
Composition is componentwise.  

However, the construction of interest is the exponential $\CC^{\CB}$ in $\mathrm{Cat}$;
it is called a {\em functor category}. We can use Example~\ref{objmorfunctors} to discover
what the objects and morphisms must be. An object of $\CC^{\CB}$ regarded as a functor $\mathbf{1}\to \CC^{\CB}$ must be obtained by currying a functor $\mathbf{1}\times \CB\to \CC$.
Yet $\mathbf{1}\times \CB$ is isomorphic to $\CB$. So the objects can be identified as functors
$S: \CB\to \CC$, hence the name ``functor category''. A morphism of $\CC^{\CB}$ regarded as a functor $\mathbf{2}\to \CC^{\CB}$ must be obtained by currying a functor 
$\mathbf{2}\times \CB\xra{R} \CC$. Such a functor $R$ gives a functor $\CB\xra{S}\CC$ and
a functor $\CB\xra{T}\CC$ with $SB=R(0,B), Sg=R(1_0,g)$ and $TB=R(1,B), Tg=R(1_1,g)$.
It also takes $(0,B)\xra{(\mu,g)}(1,D)$ to a morphism $SB\xra{\theta_g}TD$.
By applying $R$ to the commutative diagram    
\begin{eqnarray*}
\xymatrix{
(0,B) \ar[rr]^-{(\mu,1_B)} \ar[rrd]^-{(\mu,g)} \ar[d]_-{(1_0,g)} && (1,B) \ar[d]^-{(1_1,g)} \\
(0,D) \ar[rr]_-{(\mu,1_D)} && (1,D)}
\end{eqnarray*}
in $\mathbf{2}\times \CB$, we obtain the commutative diagram
\begin{eqnarray}\label{naturalitydiag}
\begin{aligned}
\xymatrix{
SB \ar[rr]^-{\theta_B} \ar[rrd]^-{\theta_g} \ar[d]_-{Sg} && TB \ar[d]^-{Tg} \\
SD \ar[rr]_-{\theta_D} && TD}
\end{aligned}
\end{eqnarray}
so that all $\theta_g$ is obtainable from the case $\theta_B=\theta_{1_B}$ where the morphism is
an identity and the domain and codomain functors $S$ and $T$.   
\begin{definition} Given functors $S, T : \CB\to \CC$, a {\em natural transformation} $\theta : S\Ra T$
from $S$ to $T$, depicted as in \eqref{natas2-cell}, consists of a morphism $\theta_B : SB\to TB$ for each object $B$ of $\CB$ such that,
for all morphisms $g : B\to D$ in $\CB$, the ``naturality condition'' $Tg\circ \theta_B = \theta_D \circ Sg$ holds. 
\begin{eqnarray}\label{natas2-cell}
 \xymatrix{
 \CB \ar@/^1pc/[rr]^S  \ar@/_1pc/[rr]_{T}  \ar@{}[rr]|{\Downarrow \; \theta}&& \CC 
 }
\end{eqnarray}  
\end{definition}

The sets $\mathrm{Cat}(\CB,\CC)$ of functors are enriched to functor categories $\CC^{\CB}$
for which the morphisms are natural transformations. This does indeed provide the exponential
for the category $\mathrm{Cat}$ where the evaluation functor 
\begin{eqnarray*}
\mathrm{ev_{\CB}} : \CC^{\CB}\times \CB \lra \CC
\end{eqnarray*}
 takes $(\theta : S\Ra T, g : B\to C)$ to the diagonal $\theta_g$ of the diagram \eqref{naturalitydiag}. 
 
 Since we have defined isomorphism in any category, this applies to the category 
 $\mathrm{Cat}$ and we can talk about isomorphic categories. 
 However, there is a weaker notion which stands in the same relation to isomorphism as isomorphism
 does to equality. A functor $T : \CA \to \CX$ is called an {\em equivalence} when there exists a 
 functor $S : \CX\to \CA$ and isomorphisms $S\circ T \cong 1_{\CA}$ and $T\circ S\cong 1_{\CX}$
 in the categories $\CA^{\CA}$ and $\CX^{\CX}$ respectively. Two categories $\CA$ and $\CX$ are
 {\em equivalent} when there exists an equivalence $\CA \to \CX$; we write $\CA\simeq \CX$. With such a sequence of relations as $=, \cong, \simeq$, you might expect that it continues. It does,
 and there is a general notion of $n$-equivalence $\sim_n$ ($n\le \infty$) in higher categories. 
 
 \begin{example}\label{Mat=Vect} The category $\mathrm{Mat}_{\mathbb{F}}$ (see Example~\ref{Matex}) is 
 actually equivalent to the category $\mathrm{vect}_{\mathbb{F}}$ of finite dimensional $\mathbb{F}$-vector spaces with linear functions as morphisms. The equivalence takes the object $n$ to the vector space $\mathbb{F}^n$.  
\end{example}
 \begin{example}\label{LinearReps} If $G$ is a finite group then $\mathrm{rep}_{\mathbb{F}}G = \mathrm{vect}_{\mathbb{F}}^G$ is the category of finite-dimensional representations of $G$
 in $\mathrm{vect}_{\mathbb{F}}$; recall Example~\ref{gpreps}. The morphisms are natural
 transformations and are traditionally called {\em $G$-equivariant linear functions}. Products in $\mathrm{rep}_{\mathbb{F}}G$ are also coproducts, are called {\em direct sums}, and are
 denoted by $\oplus$. A representation $W$ is {\em irreducible} when $W \cong U\oplus V$
 implies either $U$ or $V$ is $0$. Every representation is a direct sum of finitely many
 irreducible representations.   
\end{example}

The category $\mathrm{Cat}$ becomes what is called a {\em (strict) 2-category} when equipped with the
natural transformations as in diagram \eqref{natas2-cell}. In a 2-category, there are 
{\em 2-morphisms} between morphisms with the same domain and codomain. The concept goes
back to the early 1960s \cite{Ehres, EilKel1966}. However, $\mathrm{Cat}$ is a bit less general than
you might expect because functor composition is actually strictly associative and, as
functors are objects of a category, associativity up to isomorphism does make sense. 
Such associativity was studied by S. Mac Lane \cite{ML1963} and given the ``several object'' perspective
(recall Example~\ref{MonoidDef}) by J. B\'enabou \cite{Ben}. 
These weak 2-categories are called {\em bicategories};
see Section~\ref{Bicats}.

Even product of objects in a category is not a strictly associative operation. For example, if $X,Y,Z$ are sets, then $(X\times Y)\times Z$ is not actually equal to $X\times (Y\times Z)$: elements of the first are of the form $((x,y),z)$ while those of the second are $(x,(y,z))$. However, there is a clear bijection between such elements given
by rebracketing. In other words, we do have an isomorphism $(X\times Y)\times Z \cong X\times (Y\times Z)$. Similarly, we have isomorphisms $\mathbf{1}\times X \cong X \cong X\times \mathbf{1}$, not equalities. Of course there is also the ternary product $X\times Y\times Z$ where the elements
are $(x,y,z)$. To relate to bicategories, we look at $\mathrm{Set}$ as the hom-category of a
bicategory with a single unnamed object, with sets as morphisms, and with cartesian product
for composition.  

\section{Bicategories}\label{Bicats}

Here is the definition due to B\'enabou \cite{Ben}.
\begin{definition} A {\em bicategory} $\CW$ consists of a set $\mathrm{ob}\CW$ of objects, for each
pair of objects $U, V$, a category $\CW(U,V)$, whose objects $A : U\to V$ are called {\em morphisms} and whose morphisms $f : A\to B$ are called {\em 2-morphisms or 2-cells} (and sometimes a double
arrow $f : A\Ra B$ is used for emphasis), 
{\em identity morphisms} $I_U : U\to U$, {\em composition functors} 
\begin{eqnarray*}
\otimes : \CW(V,W)\times \CW(U,V) \lra \CW(U,W) \ ,
\end{eqnarray*}
and invertible 2-morphisms  
\begin{eqnarray}\label{iators}
 \begin{aligned}
 & a^{}_{C,B,A} : (C \otimes B) \otimes A \lra C \otimes (B \otimes A)\\
 & \ell_A : I_V \otimes A \lra A\,, \qquad r_A : A \otimes I_U \lra A
\end{aligned}
\end{eqnarray}
natural in $U\xra{A}V\xra{B}W\xra{C}X$, such that the following commute.
\begin{eqnarray}\label{pentagon}
\begin{aligned}
\footnotesize{\xymatrix{
& (D \otimes (C \otimes B)) \otimes A \ar[rd]^-{\phantom{A}a^{}_{D,C\otimes B,A}}  & \\
((D \otimes C )\otimes B) \otimes A \ar[ru]^-{a^{}_{D,C,B}\otimes 1_{A}\phantom{D}} \ar[d]_-{a^{}_{D\otimes C,B,A}} & & D \otimes ((C \otimes B) \otimes A ) \ar[d]^-{1_{D}\otimes a^{}_{C,B,A}} \\
(D \otimes C) \otimes (B \otimes A) \ar[rr]_-{a^{}_{D,C,B\otimes A }} & & D \otimes (C \otimes (B \otimes A)) 
} }
\end{aligned}
\end{eqnarray}
\begin{eqnarray}\label{unittriangle}
\begin{aligned}
\footnotesize{\xymatrix{(B \otimes I_V)\otimes A \ar[rr]^{a^{}_{B,I_V,A}} \ar[rd]_{r_B \otimes 1_A} 
& & B \otimes (I_V \otimes A) \ar[ld]^{1_B \otimes \ell_A} \\
& B \otimes A} }
\end{aligned} 
\end{eqnarray}
\end{definition}
\begin{example} A category is precisely a bicategory $\CW$ in which each of the hom-categories $\CW(U,V)$ is discrete (see Example~\ref{discretecat}).
\end{example}
\begin{example} \label{MonoidalcatDef}
If a bicategory $\CW$ has only one object $\star$ the bicategory can be identified with a {\em monoidal category} in the sense of \cite{EilKel1966}; that is, a category $\CV$ with a functorial binary operation $\CV \times \CV\xra{\otimes} \CV$ which is associative with an identity object $I$ up to isomorphisms satisfying \eqref{pentagon} and \eqref{unittriangle}. The connection is $\CV(\star,\star) = \CV$. A higher version of the slogan in Example~\ref{MonoidDef} is ``a bicategory is just a monoidal category spread over several objects''.  
\end{example}
\begin{example} Every category with finite products is monoidal by taking $\otimes = \times$.
The category $\mathrm{Mat}_{\mathbb{F}}$ whose morphisms are matrices (Example~\ref{Matex}) 
is monoidal with $m\otimes n = m+n$ since this is actual the product (and coproduct) of $m$ and $n$
in $\mathrm{Mat}_{\mathbb{F}}$. There is another monoidal structure on $\mathrm{Mat}_{\mathbb{F}}$
for which $m\otimes n = mn$ which corresponds under the equivalence of 
Example~\ref{Mat=Vect} to tensor product of vector spaces.    
\end{example}
\begin{example} A (strict) 2-category is precisely a bicategory $\CW$ in which the isomorphisms \eqref{iators} are all identities.
\end{example}
\begin{example}\label{spanbicatex} For any category $\CX$, we have the categories $\mathrm{Spn}_{\CX}(X,Y)$ for $X, Y$ objects of $\CX$ (see Section~\ref{ProdExp}).
There is a bicategory $\mathrm{Spn}_{\CX}$ where these are the hom categories. 
The composite $B\ot A$ of a span $X\xla{s_1}A\xra{s_2}Y$ and a span $Y\xla{t_1}B\xra{t_2}Z$
is $X\xla{s_1\circ \mathrm{pr}_1}A\times_{Y} B\xra{t_2\circ \mathrm{pr}_2}Z$ where $A\times_{Y} B$
is the product in the {\em slice category} $\CX_{/Y}$ whose objects are morphisms into $Y$
and whose morphisms are morphisms in $\CX$ between the domains making a commutative triangle.
In this case, we have the invertible morphisms \eqref{iators} for a bicategory which are not identities in general.   
\end{example}

Paths $U= U_0\xra{A_1} U_1 \xra{A_2} U_2\xra{A_3} \dots \xra{A_n} U_n= V$ of morphisms in a bicategory $\CW$ can be composed iteratively doing two at a time, with identities eliminated as in a category, except that (unless we have a 2-category) different bracketings lead to morphisms that are only isomorphic, not equal. However there is a unique
isomorphism made by using only the isomorphisms \eqref{iators} and the compositions $\circ$, $\otimes$. A 2-morphism between paths of morphisms means the domain and codomain of the 2-morphism are the composites of those paths. Once such composites are chosen for the outside boundary paths in so called
{\em pasting diagrams}, as in the left and right hand sides of \eqref{eqpastingdiags}, each of those diagrams determines a 2-morphism between those composites (see \cite{Ben, 7}). For example, if in
the left hand side we choose bracketing from the right for the domain and codomain paths, the composite 2-morphism is the composite of the path
\begin{eqnarray*}
G\otimes (F\otimes E)\xra{1_G\otimes g}G\ot (C'\ot A') \xra{a^{-1}_{G, C', A'}}(G\ot C')\ot A' \xra{h\ot 1_{A'}} \\
(D\ot B')\ot A' \xra{a_{D,B',A'}} D\ot (B'\ot A') \xra{1_D\ot f} D\ot(C\ot (B\ot A))
\end{eqnarray*}
in the category $\CW(U,X)$.

\begin{eqnarray}\label{eqpastingdiags}
\begin{aligned}
\xymatrix{\ar @{} [rdd] | {\stackrel{g} \Longrightarrow}
U \ar[r]^-{A} \ar[rd]|-{A'} \ar[dd]_-{E} &  \ar @{} [rd] | {\stackrel{f} \Longrightarrow} P \ar[r]^-{B} &  V \ar[d]^{C} \\
& M \ar[r]|-{B'} \ar[d]|-{C'} & R \ar[d]^{D} \\
W \ar[r]_{F} & Q \ar[r]_{G} & X  \ar @{} [lu] | {\stackrel{h} \Longrightarrow}
}
\end{aligned}
\qquad
=
\qquad
\begin{aligned}
\xymatrix{\ar @{} [rd] |  {\stackrel{m} \Longrightarrow}
U \ar[r]^-{A}  \ar[dd]_-{E} &  P\ar @{} [rd] | {\stackrel{\ell} \Longrightarrow}  \ar[d]|-{D'} \ar[rr]^-{B} & &  V  \ar[ld]|-{E'} \ar[d]^{C} \\
& N \ar[r]|-{F'}  & L \ar[dr]|-{G'}  & R \ar[d]^{D} \\
W \ar[ru]|-{H} \ar[r]_{F} & Q \ar[rr]_{G} & & X \ar @{} [lllu] | {\Uparrow k} \ar @{} [luu] | {\stackrel{n} \Longrightarrow}   
}
\end{aligned}
\end{eqnarray}
 
\section{Strings}
From the outset, categories were to connect algebra and geometry.
Both algebraic structures and geometric structures based on sets had their morphisms
which were functions with appropriate properties. Prime examples of functors were those
constructing algebraic structures from geometric ones or vice versa. The connection I wish
to discuss now is somewhat different.

Monoids are sets with algebraic structure. Categories are graphs with algebraic structure.
However, there is geometry inherent in the notation we use in a category: the diagrams involve
0-dimensional objects, 1-dimensional arrows denoting morphisms, and 2-dimensional regions in diagrams. In a bicategory there is also a 3-dimensional aspect: notice that, 
in diagram \eqref{eqpastingdiags}, the outside boundary paths of the two sides are equal and we
can glue them over each other creating a bubble (or 2-globe) with the equality between the diagrams
becoming the 3-dimensional interior. The reason we have drawn \eqref{eqpastingdiags} the way we
have, repeating the boundaries, is because paper is 2-dimensional. While there is this precise 
geometry in pasting diagrams (see \cite{Pow}), this is not the way monoidal categories
 and bicategories made their recent impact on the theory of manifolds and theoretical physics.

Many subjects have traditionally involved connecting things with wires or strings to understand
how the things are related and even to perform calculations; examples are circuit diagrams, networks, Feynman diagrams in relativity, Petri nets, flow charts, and planar projections in knot theory.
Our strings can be seen as adapting the use Roger Penrose \cite{Pen1971} made of them to connect 
repeated indices on the tensors (like $\tensor{R \ }{^\mu_\nu^\rho_\tau^\sigma}$) coming from differential geometry and relativity. 

The starting point is to replace the globular diagram
\begin{eqnarray}\label{2-cell}
\footnotesize{ \xymatrix{
U \ar@/^1pc/[rr]^A  \ar@/_1pc/[rr]_{B}  \ar@{}[rr]|{\Downarrow \; f}&& V 
 }}
\end{eqnarray}  
for a 2-morphism in a bicategory $\CW$ by a string diagram
\begin{eqnarray}
\begin{aligned}
\psscalebox{0.6 0.6} % Change this value to rescale the drawing.
{
\begin{pspicture}(0,-1.6803068)(4.58,1.6803068)
\psellipse[linecolor=black, linewidth=0.04, dimen=outer](2.41,-0.2699387)(0.33,0.33)
\psline[linecolor=black, linewidth=0.04](2.4,-0.6)(2.4,-1.7)
\rput[bl](4.2,-0.3999387){$V$}
\rput[bl](2.0,-1.7){\scriptsize{$B$}}
\rput[bl](2.0,1.1800613){\scriptsize{$A$}}
\rput[bl](2.36,-0.3999387){\scriptsize{$f$}}
\rput[bl](0.0,-0.3999387){$U$}
\psline[linecolor=black, linewidth=0.04](2.4,0.07)(2.4,1.4)
\end{pspicture}
}
\end{aligned}
\end{eqnarray}
obtained by a dimension swap: 0-dimensional objects are changed to 2-dimensional regions.
A composite $A\xra{f}B\xra{g}C$ of 2-morphisms between morphisms $A, B, C$ from $U$ to $V$ is obtained by placing
$g$ below $f$ and identifying the two strings labelled $B$. The operation $\circ$ is thereby built into the geometry and is often called {\em vertical composition}. An identity 2-morphism $1_A : A\Ra A$ where $A : U\to V$ is depicted by a single vertical string labelled $A$
with region $U$ to the left and $V$ to the right. The operation $\otimes$ is also built into the geometry and is often called {\em horizontal composition}; it is performed
on string diagrams by placing the component diagrams horizontally next to each other.
An identity morphism $I_U : U\to U$ is depicted as an unlabelled vertical string with region
$U$ to the left and right.
A pasting diagram, or component thereof, as in \eqref{pastingCAfEDB} where $f : C\ot A\Ra E\ot D\ot B$,
\begin{equation}\label{pastingCAfEDB}
\begin{aligned}
\footnotesize{\xymatrix{
U \ar[d]_{A}^(0.8){\phantom{aaaa}}="1" \ar[rr]^{B}  && Y \ar[d]^{D}_(0.8){\phantom{aaaa}}="2" \ar@{=>}"1";"2"^-{f}
\\
X \ar[rd]_-{C} && Z \ar[ld]^-{E} 
\\
& V 
}}
\end{aligned}
\end{equation}
becomes the string diagram \eqref{stringCAfEDB}.
\begin{equation}\label{stringCAfEDB}
\begin{aligned}
\psscalebox{0.6 0.6} % Change this value to rescale the drawing.
{
\begin{pspicture}(0,-1.5917871)(8.16,1.5917871)
\psellipse[linecolor=black, linewidth=0.04, dimen=outer](5.13,-0.06584961)(0.33,0.33)
\psline[linecolor=black, linewidth=0.04](4.88,0.08415039)(2.26,1.5641503)
\psline[linecolor=black, linewidth=0.04](5.4,0.12415039)(7.44,1.5641503)
\psline[linecolor=black, linewidth=0.04](4.86,-0.2358496)(2.28,-1.4158496)
\psline[linecolor=black, linewidth=0.04](5.38,-0.2558496)(7.7,-1.3758496)
\psline[linecolor=black, linewidth=0.04](5.12,-0.39584962)(5.1,-1.4758496)
\rput[bl](5.05,-0.21584961){\scriptsize{$f$}}
\rput[bl](7.75,-1.30){\scriptsize{$E$}}
\rput[bl](2.2,-1.2558496){\scriptsize{$B$}}
\rput[bl](7.1,1.1041504){\scriptsize{$C$}}
\rput[bl](5.2,-1.5158496){\scriptsize{$D$}}
\rput[bl](2.84,-0.055849608){$U$}
\rput[bl](6.94,-0.07584961){$V$}
\rput[bl](4.9,1.3641504){$X$}
\rput[bl](3.74,-1.2958496){$Y$}
\rput[bl](5.88,-1.2958496){$Z$}
\rput[bl](2.3,1.1041504){\scriptsize{$A$}}
\end{pspicture}
}
\end{aligned}
\end{equation}
An equation between pasting diagrams such as \eqref{eqpastingdiags} becomes the equation
between string diagrams as in \eqref{eqstringdiags}.  
\begin{equation}\label{eqstringdiags}
\begin{aligned}
\psscalebox{0.5 0.5} % Change this value to rescale the drawing.
{
\begin{pspicture}(0,-3.3822997)(20.38,3.3822997)
\pscircle[linecolor=black, linewidth=0.04, dimen=outer](2.04,1.6736379){0.32}
\pscircle[linecolor=black, linewidth=0.04, dimen=outer](5.0,-0.106362194){0.32}
\pscircle[linecolor=black, linewidth=0.04, dimen=outer](2.0,-1.7263622){0.32}
\pscircle[linecolor=black, linewidth=0.04, dimen=outer](14.96,2.153638){0.32}
\pscircle[linecolor=black, linewidth=0.04, dimen=outer](12.92,0.9936378){0.32}
\pscircle[linecolor=black, linewidth=0.04, dimen=outer](14.4,-0.5663622){0.32}
\pscircle[linecolor=black, linewidth=0.04, dimen=outer](16.42,-1.7263622){0.32}
\rput[bl](1.96,1.5536379){\scriptsize{$g$}}
\rput[bl](4.94,-0.22636218){\scriptsize{$h$}}
\rput[bl](1.92,-1.8463622){\scriptsize{$f$}}
\rput[bl](12.84,0.8736378){\scriptsize{$m$}}
\rput[bl](16.36,-1.8463622){\scriptsize{$n$}}
\rput[bl](14.88,2.0336378){\scriptsize{$k$}}
\rput[bl](14.32,-0.6863622){\scriptsize{$\ell$}}
\psline[linecolor=black, linewidth=0.04](1.84,1.8736378)(1.12,3.373638)
\psline[linecolor=black, linewidth=0.04](2.24,1.8936378)(3.14,3.3336377)
\psline[linecolor=black, linewidth=0.04](5.14,0.1736378)(6.28,3.2536378)
\psline[linecolor=black, linewidth=0.04](2.28,1.4336379)(4.78,0.11363781)
\psline[linecolor=black, linewidth=0.04](1.84,1.4936378)(1.84,-1.4863622)
\psline[linecolor=black, linewidth=0.04](2.24,-1.5263622)(4.82,-0.3463622)
\psline[linecolor=black, linewidth=0.04](5.18,-0.3263622)(6.24,-3.346362)
\psline[linecolor=black, linewidth=0.04](1.76,-1.9263622)(0.76,-3.3063622)
\psline[linecolor=black, linewidth=0.04](2.04,-2.0263622)(2.08,-3.346362)
\psline[linecolor=black, linewidth=0.04](2.26,-1.9263622)(3.9,-3.3263621)(3.92,-3.346362)
\psline[linecolor=black, linewidth=0.04](12.7,1.1936378)(11.4,3.2736378)
\psline[linecolor=black, linewidth=0.04](13.1,1.2336378)(14.72,1.9536378)
\psline[linecolor=black, linewidth=0.04](14.9,1.8536378)(14.6,-0.3463622)
\psline[linecolor=black, linewidth=0.04](13.1,0.7336378)(14.24,-0.36636218)
\psline[linecolor=black, linewidth=0.04](12.76,0.7136378)(11.76,-3.3063622)
\psline[linecolor=black, linewidth=0.04](14.18,-0.7863622)(13.08,-3.3263621)
\psline[linecolor=black, linewidth=0.04](14.64,-0.8063622)(16.18,-1.5463622)
\psline[linecolor=black, linewidth=0.04](15.18,1.9336379)(16.62,-1.5063622)
\psline[linecolor=black, linewidth=0.04](16.2,-1.9463621)(15.48,-3.3063622)
\psline[linecolor=black, linewidth=0.04](16.66,-1.9463621)(17.9,-3.3263621)
\psline[linecolor=black, linewidth=0.04](14.74,2.3336377)(14.04,3.2536378)
\psline[linecolor=black, linewidth=0.04](15.16,2.373638)(16.2,3.2536378)
\rput[bl](8.9,-0.12636219){\Large{$=$}}
\rput[bl](0.0,0.15363781){\large{$U$}}
\rput[bl](2.34,-2.846362){\large{$V$}}
\rput[bl](1.66,2.9536378){\large{$W$}}
\rput[bl](6.2,0.11363781){\large{$X$}}
\rput[bl](1.38,-2.866362){\large{$P$}}
\rput[bl](3.8,1.6){\large{$Q$}}
\rput[bl](14.96,0.01363781){\large{$L$}}
\rput[bl](13.58,0.77363783){\large{$N$}}
\rput[bl](4.34,-2.5263622){\large{$R$}}
\rput[bl](16.22,-2.8263621){\large{$R$}}
\rput[bl](17.0,0.13363782){\large{$X$}}
\rput[bl](10.78,0.09363781){\large{$U$}}
\rput[bl](12.72,-1.2463622){\large{$P$}}
\rput[bl](14.36,-2.386362){\large{$V$}}
\rput[bl](12.62,2.8336377){\large{$W$}}
\rput[bl](14.68,2.8536377){\large{$Q$}}
\rput[bl](2.58,-0.106362194){\large{$M$}}
\rput[bl](12.8,-3.2463622){\footnotesize{$B$}}
\rput[bl](1.72,-3.3063622){\footnotesize{$B$}}
\rput[bl](15.2,-3.2263622){\footnotesize{$C$}}
\rput[bl](3.3,-3.3063622){\footnotesize{$C$}}
\rput[bl](17.0,-2.9263623){\footnotesize{$D$}}
\rput[bl](5.7,-2.9863622){\footnotesize{$D$}}
\rput[bl](0.8,3.0936377){\footnotesize{$E$}}
\rput[bl](11.1,2.9336379){\footnotesize{$E$}}
\rput[bl](13.7,2.9536378){\footnotesize{$F$}}
\rput[bl](2.6,2.9336379){\footnotesize{$F$}}
\rput[bl](15.65,3.0536377){\footnotesize{$G$}}
\rput[bl](5.76,2.9136379){\footnotesize{$G$}}
\rput[bl](13.32,1.5336378){\footnotesize{$H$}}
\rput[bl](1.4,-0.40636218){\footnotesize{$A'$}}
\rput[bl](3.12,-0.9863622){\footnotesize{$B'$}}
\rput[bl](0.48,-3.2063622){\footnotesize{$A$}}
\rput[bl](11.48,-2.886362){\footnotesize{$A$}}
\rput[bl](13.18,0.03363781){\footnotesize{$D'$}}
\rput[bl](14.4,0.41363782){\footnotesize{$F'$}}
\rput[bl](2.78,0.6936378){\footnotesize{$C'$}}
\rput[bl](15.64,-0.30636218){\footnotesize{$G'$}}
\rput[bl](14.86,-1.3463622){\footnotesize{$E'$}}
\end{pspicture}
}
\end{aligned}
\end{equation}

When dealing with a monoidal category (see Example~\ref{MonoidalcatDef}), we use the term {\em tensor product} for the horizontal composition $\ot$.
In our paper \cite{GTC}, we proved that the string diagrams in the plane could be used
for valid calculations in monoidal categories.

\section{Adjunction}\label{Adj}
Here is a fundamental categorical concept which was formulated for functors by Daniel Kan
\cite{Kan1958}. 
\begin{definition}\label{adjunctiondef} A morphism $V\xra{B} U$ is said to be {\em left adjoint} to a morphism
$U\xra{A} V$ in a bicategory $\CW$ when there exist 2-morphisms $e : B\ot A\Ra 1_U$
and $d : 1_V\Ra A\ot B$ such that the two equations \eqref{snakes} hold. The 2-morphisms
$d$ and $e$ are called the {\em unit} and {\em counit}. When they exist
the notation $B\dashv A$ is used. If $B\dashv A$ and $B'\dashv A$, then there is
an invertible 2-morphism $B\cong B'$ compatible with the units and counits. 
We also say $A$ is {\em right adjoint} to $B$ when $B\dashv A$.
\end{definition}
\begin{eqnarray}\label{snakes}
\begin{aligned}
\psscalebox{0.5 0.5} % Change this value to rescale the drawing.
{
\begin{pspicture}(0,-3.2428024)(19.94,3.2428024)
\pscircle[linecolor=black, linewidth=0.04, dimen=outer](1.38,1.1573204){0.36}
\pscircle[linecolor=black, linewidth=0.04, dimen=outer](3.54,-0.7226797){0.36}
\rput[bl](1.32,1.0373204){\footnotesize{$d$}}
\rput[bl](3.48,-0.8426797){\footnotesize{$e$}}
\pscircle[linecolor=black, linewidth=0.04, dimen=outer](14.92,1.0773203){0.36}
\rput[bl](14.86,0.95732033){\footnotesize{$d$}}
\pscircle[linecolor=black, linewidth=0.04, dimen=outer](12.76,-0.74267966){0.36}
\rput[bl](12.7,-0.86267966){\footnotesize{$e$}}
\psline[linecolor=black, linewidth=0.04](1.14,0.89732033)(0.22,-3.0626798)
\psline[linecolor=black, linewidth=0.04](1.6,0.9173204)(3.28,-0.50267965)
\psline[linecolor=black, linewidth=0.04](3.8,-0.48267967)(4.86,3.2373204)
\psline[linecolor=black, linewidth=0.04](12.5,-0.5226796)(11.1,3.2173204)
\psline[linecolor=black, linewidth=0.04](12.98,-0.48267967)(14.66,0.85732037)
\psline[linecolor=black, linewidth=0.04](15.14,0.79732037)(16.2,-3.0826797)
\rput[bl](0.0,-1.4826796){\small{$A$}}
\rput[bl](1.92,-0.042679653){\small{$B$}}
\rput[bl](3.88,2.0173204){\small{$A$}}
\rput[bl](10.9,2.0173204){\small{$B$}}
\rput[bl](15.1,-1.3226796){\small{$B$}}
\rput[bl](13.22,0.117320344){\small{$A$}}
\rput[bl](5.0,0.41732034){\LARGE{$=$}}
\psline[linecolor=black, linewidth=0.04](6.66,3.1973205)(6.62,-3.1826797)
\rput[bl](16.5,0.43732035){\LARGE{$=$}}
\psline[linecolor=black, linewidth=0.04](18.64,3.1373203)(18.6,-3.2426796)
\rput[bl](2.24,2.2173204){$U$}
\rput[bl](14.04,-2.2226796){$V$}
\rput[bl](13.76,2.3373203){$U$}
\rput[bl](7.22,-2.2026796){$V$}
\rput[bl](1.9,-2.2426796){$V$}
\rput[bl](6.06,2.2773204){\small{$A$}}
\rput[bl](18.1,2.3373203){\small{$B$}}
\rput[bl](17.66,-2.3626797){$V$}
\rput[bl](19.3,1.3773203){$U$}
\rput[bl](5.78,1.3973204){$U$}
\end{pspicture}
}
\end{aligned}
\end{eqnarray}
\begin{example} The functor $\mathrm{Mon}\xra{T} \mathrm{Set}$, which takes each monoid
to the set underlying the monoid, has a left adjoint $\mathrm{Set}\xra{S} \mathrm{Mon}$ 
(in the 2-category of categories) taking
each set $\Lambda$ to the monoid $S\Lambda = \Lambda^*$ of words in the alphabet
$\Lambda$; see Example~\ref{alphabet}. 
\end{example}

\begin{example}\label{monoidaldual} If the bicategory has one object and so amounts to a 
monoidal category $\CV$ then we say $B$ is a {\em left (monoidal) dual of $A$}
when $B\dashv A$ and we write $B = A^{\vee}$ or $A = B^{\wedge}$. 
If $\CV = \mathrm{vect}_{\mathbb{F}}$ (see Example~\ref{Mat=Vect}) 
then the vector space of linear functions from $A$ to $\mathbb{F}$ is a 
left dual $A^{\vee}$ for $A$ with the counit $e : A^{\vee}\ot A\to \mathbb{F}$
taken to be the evaluation function $e(\phi \ot x) = \phi(x)$. In fact, a vector space
in the symmetric monoidal category of all vector spaces has a dual if and only if
it is finite dimensional.    
\end{example}

A monoidal category $\CV$ is called {\em autonomous} when every object has both a left and right dual. If $\CV$ is braided, the left duals are also right duals and $\CV$ is called {\em rigid}
and called {\em compact} if the braiding is a symmetry. 

\section{Enriched categories}\label{enriched}

For an ordinary category $\CA$, the homs $\CA(A,B)$ are sets whose elements are the
morphisms from $A$ to $B$. It is often the case that these sets have more structure.
For example, if $\CA$ is the category of vector spaces over a field, each hom $\CA(A,B)$
is a vector space and composition \eqref{compo} is bilinear. Bilinearity here means that
composition is defined on a tensor product of the homs as a linear function, not just a function on the cartesian product. This suggests generalising the notion of category to one where
the homs lie in some monoidal category $\CV$ (recall Example~\ref{MonoidalcatDef}). The objects do not even need to be sets with structure. We are thus led to the definition of a {\em category with homs enriched in the base $\CV$} and briefly called a {\em $\CV$-enriched category}.  

A monoidal category $\CV$ is called {\em right closed} when, for each object $X$, the functor
$X\ot - : \CV\to \CV$ has a right adjoint. The value $Y^X$ of the right adjoint at $Y$ is
called the {\em internal hom} of $X, Y$ for $\CV$. There is a natural bijection
\begin{eqnarray*}
\CV(X\ot Z, Y) \cong \CV(Z, Y^X) \ . 
\end{eqnarray*}
In this case $\CV$ becomes a $\CV$-enriched category with $Y^X$ as $\CV$-valued hom. When the tensor product is product, $Y^X$ is the exponential as in Section~\ref{ProdExp}. 

For example, a 2-category is a $\mathrm{Cat}$-category; see \cite{EilKel1966, KelBk}.
Inductively we can define an {$(n+1)$-category} to be an $n$-Cat-enriched category
where the tensor product in each category $n$-Cat is product. It is harder to 
precisely define weak $n$-categories in the sense I used it in \cite{30} (with a hint on
an approach perfected by Dominic Verity \cite{Ver}) but a number of authors have 
different definitions. Comparison of the different approaches is an important ongoing study.

A Banach space is a vector space $B$ over the field $\mathbb{C}$ of complex numbers equipped with a function $\lvert-\rvert : B\to \mathbb{R}$, called the {\em norm}, whose value $\lvert x\rvert$
at $x$ in $B$ is to be regarded as the distance from $x$ to the origin $0$; there are various
axioms. There is a monoidal category $\mathrm{Ban}$ of Banach spaces and norm-non-increasing linear functions. The internal hom $Y^X$ is the space of bounded linear functions from $X$ to $Y$.    

\section{Quantum theory}\label{Qt}

Through the work of Hermann Weyl \cite{Wey1918, Wey1927} which aimed to provide 
mathematical foundations for general relativity and quantum mechanics, groups and their representations came into focus. The groups of symmetries of the systems 
(now not just finite as in Example~\ref{gpreps}) have extra structure as smooth manifolds
 and the representations are required to respect that structure. 
In quantum mechanics, the states can be labelled by irreducible representations of the group of symmetries. 
In quantum field theory, the states amount to particles. 

The group of symmetries is often a compact group $G$ such as the group of $n \times n$-matrices with complex entries and with the conjugate transpose as inverse.
The representations of interest are unitary and are good examples of enriched functors
between manifold-enriched categories. 
The category $\mathrm{Rep}G$ of such representations is a compact monoidal $\mathrm{Ban}$-enriched category. (This example was the reason ``compact''
was suggested by Peter Freyd for the monoidal category context.) 
 
 In 1971 Freyd was visiting Max Kelly at the University of NSW. He gave seminars on
 a question asked by the mathematical physicist John Roberts working on superselection principles in quantum field theory. Roberts thought it quite possible that $\mathrm{Rep}G$ might be observable as a monoidal category before knowing $G$. 
 The question was: given a symmetric monoidal $C^*$-category $\CC$, 
 is there a compact group $G$ for which $\CC$ is equivalent $\mathrm{Rep}G$? 
 A $C^*$-category is a $C^*$-algebra ``with several objects'' and $C^*$-algebras 
 were already well used in mathematics (functional analysis) and quantum theory.
Finite groups are compact and Freyd dealt with that case in his seminars. 
The full answer had to wait until 1989 to be published by Doplicher-Roberts \cite{DopRob}.
 
 Already in 1964 Alexander Grothendieck had constructed a category of ``motives''
 with a structure preserving functor into a category of vector spaces which led him
 to the so-called Grothendieck-Galois (pro-algebraic) group. He conjectured that the representations of the group gave back the motives. Following this came the development of classical Tannaka duality \cite{Tan, Che} in the setting of algebraic geometry and category theory \cite{Saa, DelMil}.  
 
 Later in the 1970s, Roberts visited Australia a couple of times. He talked in our category seminar on net cohomology and quantum field theory. The idea that appealed to me was that he thought this non-abelian cohomology should have algebraic coefficients which were higher-dimensional categories: $n$-categories where $n$ could be $\omega$, the first infinite ordinal. 
For example, the associativity pentagon \eqref{pentagon} is actually the cohomological 3-cocycle condition.  
Moreover, so that the topology could relate to the algebra, he wanted to characterize $n$-categories as
 ``simplicial sets'' which are combinatorial structures like graphs that had been studied by
 Eilenberg-Zilber and Kan in connection with functors, called homology and homotopy,
 that provide algebraic invariants for geometric structures such as topological spaces. 
 
 Roberts defined ``complicial sets'' which were simplicial sets with some special elements
 satisfying some unique ``horn filler'' existence conditions. He expected the category of these to be equivalent to the category of $\omega$-categories. This required a functor going each way 
 between the categories of simplicial sets and the category of $\omega$-categories. 
 I realised that I could find an adjoint pair of functors (see Section~\ref{Adj}) if I could construct
 the free $\omega$-category $\mathbf{O}_n$ on the $n$-simplex which I called the {\em $n$-th oriental}. 
 I found a model for the $\omega$-category $\mathbf{O}_n$ in 1982 and later proved its freeness property;
 see \cite{30}. When restricted to complicial sets, the adjunction gave the conjectured equivalence, as proved by Verity \cite{VerO}.              

\section{Braids, knots, links and tangles}\label{Bklt}

When the monoidal category $\CV$ is equipped with the extra structure, called a {\em braiding} in \cite{ix, BTC}, our string diagrams move into three dimensions and our paper \cite{GTC} also proved that these 3-dimensional diagrams could be used for valid calculations in these braided monoidal categories.

A braiding on a monoidal category $\CV$ is a family of invertible morphisms 
\begin{eqnarray*}
c_{A,B} : A \otimes B  \lra  B \otimes A
\end{eqnarray*}
satisfying four equations: the two equations \eqref{braidaxx} and their rotations through $180^{\circ}$, for all $A\xra{f}A'$.
\begin{equation}\label{braidaxx}
\begin{aligned}
\psscalebox{0.51 0.51} % Change this value to rescale the drawing.
{
\begin{pspicture}(0,-2.4898353)(23.46,2.4898353)
\pscircle[linecolor=black, linewidth=0.04, dimen=outer](2.1,0.040760994){0.32}
\rput[bl](0.0,1.860761){\small{$A$}}
\rput[bl](2.56,-1.919239){\small{$A$}}
\psline[linecolor=black, linewidth=0.04](1.84,0.14076099)(0.44,2.380761)
\psline[linecolor=black, linewidth=0.04](1.96,0.340761)(0.66,2.460761)
\psline[linecolor=black, linewidth=0.04](3.44,-2.479239)(2.04,-0.239239)
\psline[linecolor=black, linewidth=0.04](3.56,-2.359239)(2.26,-0.21923901)
\psline[linecolor=black, linewidth=0.04](2.26,0.30076098)(3.42,2.420761)
\psline[linecolor=black, linewidth=0.04](1.92,-0.199239)(0.52,-2.319239)
\rput[bl](1.08,1.840761){\small{$B$}}
\rput[bl](3.36,-1.939239){\small{$B$}}
\rput[bl](1.98,-0.039239004){\footnotesize{$c$}}
\rput[bl](1.02,-1.919239){\small{$C$}}
\rput[bl](3.3,1.880761){\small{$C$}}
\rput[bl](4.74,-0.09923901){\LARGE{$=$}}
\pscircle[linecolor=black, linewidth=0.04, dimen=outer](7.72,-0.719239){0.28}
\pscircle[linecolor=black, linewidth=0.04, dimen=outer](9.18,1.320761){0.28}
\psline[linecolor=black, linewidth=0.04](9.06,1.080761)(7.9,-0.539239)
\psline[linecolor=black, linewidth=0.04](8.96,1.460761)(8.44,2.480761)
\psline[linecolor=black, linewidth=0.04](9.4,1.460761)(9.92,2.460761)
\psline[linecolor=black, linewidth=0.04](7.56,-0.939239)(6.66,-2.239239)
\psline[linecolor=black, linewidth=0.04](7.9,-0.959239)(8.58,-2.239239)
\psline[linecolor=black, linewidth=0.04](7.54,-0.539239)(7.12,2.460761)
\psline[linecolor=black, linewidth=0.04](9.34,1.080761)(9.94,-2.239239)
\rput[bl](9.06,1.240761){\footnotesize{$c$}}
\rput[bl](7.62,-0.799239){\footnotesize{$c$}}
\rput[bl](6.74,2.000761){\small{$A$}}
\rput[bl](7.78,-1.859239){\small{$A$}}
\rput[bl](8.12,2.020761){\small{$B$}}
\rput[bl](9.96,-1.839239){\small{$B$}}
\rput[bl](8.1,0.260761){\small{$C$}}
\rput[bl](6.48,-1.879239){\small{$C$}}
\rput[bl](9.92,2.060761){\small{$C$}}
\pscircle[linecolor=black, linewidth=0.04, dimen=outer](15.56,-0.659239){0.28}
\rput[bl](12.06,0.20076099){\LARGE{$,$}}
\pscircle[linecolor=black, linewidth=0.04, dimen=outer](19.66,1.120761){0.28}
\rput[bl](17.5,-0.039239004){\LARGE{$=$}}
\pscircle[linecolor=black, linewidth=0.04, dimen=outer](14.82,1.140761){0.28}
\pscircle[linecolor=black, linewidth=0.04, dimen=outer](20.4,-0.61923903){0.28}
\psline[linecolor=black, linewidth=0.04](14.74,1.340761)(14.34,2.420761)
\psline[linecolor=black, linewidth=0.04](14.92,0.900761)(15.48,-0.419239)
\psline[linecolor=black, linewidth=0.04](15.7,-0.439239)(16.34,2.460761)
\psline[linecolor=black, linewidth=0.04](15.44,-0.919239)(14.88,-2.219239)(14.88,-2.259239)
\psline[linecolor=black, linewidth=0.04](15.72,-0.899239)(16.22,-2.259239)
\psline[linecolor=black, linewidth=0.04](19.48,1.340761)(18.78,2.440761)
\psline[linecolor=black, linewidth=0.04](19.82,1.320761)(20.38,2.460761)
\psline[linecolor=black, linewidth=0.04](19.78,0.860761)(20.28,-0.399239)
\psline[linecolor=black, linewidth=0.04](20.48,-0.879239)(20.78,-2.159239)
\psline[linecolor=black, linewidth=0.04](19.52,0.900761)(19.08,-2.199239)
\rput[bl](15.44,-0.779239){\footnotesize{$c$}}
\rput[bl](19.54,1.000761){\footnotesize{$c$}}
\rput[bl](14.7,1.040761){\footnotesize{$f$}}
\rput[bl](20.28,-0.739239){\footnotesize{$f$}}
\rput[bl](13.88,2.120761){\small{$A$}}
\rput[bl](18.48,2.120761){\small{$A$}}
\rput[bl](20.12,0.160761){\small{$A$}}
\rput[bl](16.36,2.100761){\small{$B$}}
\rput[bl](20.46,2.140761){\small{$B$}}
\rput[bl](14.6,-1.8992391){\small{$B$}}
\rput[bl](18.66,-1.879239){\small{$B$}}
\rput[bl](20.84,-1.819239){\small{$A'$}}
\rput[bl](14.66,0.240761){\small{$A'$}}
\rput[bl](16.14,-1.859239){\small{$A'$}}
\end{pspicture}
}
\end{aligned}
\end{equation}
A braiding is called a {\em symmetry} when $c_{B,A}\circ c_{A,B} = 1_{A\ot B}$;
a monoidal category with such is called {\em symmetric}.

The way 3-dimensional space enters the picture is by realising that equations \eqref{braidaxx}
are geometrically true if we replace the nodes labelled $c$ by a string crossing in three
dimensions: left string over right as in the left diagram of \eqref{crossing/YB}.
The inverse of the braiding is then depicted by a right over left crossing.
The equation of \eqref{crossing/YB} can be proved algebraically from the axioms \eqref{braidaxx} but is clearly true geometrically; it is a Coxeter relation for the braid group,
a Reidemeister move in knot theory, and the Yang-Baxter equation in quantum theory.  

\begin{equation}\label{crossing/YB}
\begin{aligned}
\psscalebox{0.7 0.7} % Change this value to rescale the drawing.
{
\begin{pspicture}(0,-1.462688)(2.5154603,1.462688)
\psline[linecolor=black, linewidth=0.04](0.12,1.4500002)(2.5,-1.4499998)
\psline[linecolor=black, linewidth=0.04](2.48,1.4300002)(1.42,0.07000015)
\psline[linecolor=black, linewidth=0.04](1.2,-0.12999985)(0.12,-1.4099998)
\rput[bl](0.0,0.7500002){\small{$A$}}
\rput[bl](0.0,-0.8699998){\small{$B$}}
\end{pspicture}
}
\end{aligned}
\qquad
,
\qquad
\begin{aligned}
\psscalebox{0.55 0.55} % Change this value to rescale the drawing.
{
\begin{pspicture}(0,-2.3233075)(9.811152,2.3233075)
\psline[linecolor=black, linewidth=0.04](0.29557633,2.3108344)(3.9555764,-2.2691658)
\psline[linecolor=black, linewidth=0.04](3.8755763,2.2708342)(2.2555764,0.25083432)
\psline[linecolor=black, linewidth=0.04](1.8755764,-0.069165684)(1.3155763,-0.7291657)
\psline[linecolor=black, linewidth=0.04](1.0155764,-1.0091656)(0.015576324,-2.2491658)
\psline[linecolor=black, linewidth=0.04](2.0555763,2.2508342)(1.2955763,1.4108343)
\pscustom[linecolor=black, linewidth=0.04]
{
\newpath
\moveto(0.97557634,1.1708343)
\lineto(0.9255763,1.1108344)
\curveto(0.90057635,1.0808343)(0.85057634,1.0208343)(0.8255763,0.9908343)
\curveto(0.8005763,0.9608343)(0.7505763,0.9008343)(0.72557634,0.8708343)
\curveto(0.7005763,0.8408343)(0.67057633,0.6108343)(0.66557634,0.4108343)
\curveto(0.66057634,0.21083431)(0.6805763,-0.054165687)(0.7055763,-0.11916569)
\curveto(0.73057634,-0.18416569)(0.78057635,-0.3091657)(0.8055763,-0.3691657)
\curveto(0.8305763,-0.4291657)(0.8805763,-0.5291657)(0.90557635,-0.5691657)
\curveto(0.9305763,-0.60916567)(0.98057634,-0.67916566)(1.0055764,-0.7091657)
\curveto(1.0305763,-0.73916566)(1.0805763,-0.80416566)(1.1055763,-0.8391657)
\curveto(1.1305764,-0.8741657)(1.1805763,-0.9491657)(1.2055763,-0.98916566)
\curveto(1.2305763,-1.0291657)(1.2805763,-1.1091657)(1.3055763,-1.1491656)
\curveto(1.3305763,-1.1891657)(1.3805764,-1.2691656)(1.4055763,-1.3091657)
\curveto(1.4305763,-1.3491657)(1.4805763,-1.4191657)(1.5055764,-1.4491657)
\curveto(1.5305763,-1.4791657)(1.5805763,-1.5341657)(1.6055763,-1.5591657)
\curveto(1.6305764,-1.5841657)(1.6855763,-1.6341656)(1.7155763,-1.6591657)
\curveto(1.7455764,-1.6841657)(1.8155763,-1.7591656)(1.8555763,-1.8091657)
\curveto(1.8955764,-1.8591657)(1.9605763,-1.9541657)(1.9855763,-1.9991657)
\curveto(2.0105762,-2.0441656)(2.0505762,-2.1391656)(2.0655763,-2.1891656)
}
\psline[linecolor=black, linewidth=0.04](6.0355763,2.2708342)(9.695577,-2.3091657)
\psline[linecolor=black, linewidth=0.04](7.735576,-0.18916568)(6.0555763,-2.2691658)
\psline[linecolor=black, linewidth=0.04](9.795576,2.2308342)(8.795576,0.9908343)
\psline[linecolor=black, linewidth=0.04](8.535576,0.7308343)(8.055576,0.17083432)
\pscustom[linecolor=black, linewidth=0.04]
{
\newpath
\moveto(7.8555765,2.2308342)
\lineto(7.905576,2.1608343)
\curveto(7.9305763,2.1258342)(7.9805765,2.0608344)(8.005576,2.0308342)
\curveto(8.030577,2.0008342)(8.080576,1.9258343)(8.1055765,1.8808343)
\curveto(8.130576,1.8358343)(8.180576,1.7358344)(8.205576,1.6808343)
\curveto(8.2305765,1.6258343)(8.280577,1.5258343)(8.305576,1.4808344)
\curveto(8.330576,1.4358343)(8.380576,1.3608344)(8.405577,1.3308343)
\curveto(8.430576,1.3008343)(8.4805765,1.2158343)(8.505576,1.1608343)
\curveto(8.530577,1.1058344)(8.580576,1.0108343)(8.6055765,0.9708343)
\curveto(8.630576,0.9308343)(8.680576,0.8508343)(8.705576,0.8108343)
\curveto(8.7305765,0.7708343)(8.780577,0.6908343)(8.805576,0.6508343)
\curveto(8.830576,0.6108343)(8.880576,0.5408343)(8.905577,0.51083434)
\curveto(8.930576,0.4808343)(8.9805765,0.4108343)(9.005576,0.37083432)
\curveto(9.030577,0.3308343)(9.080576,0.22583431)(9.1055765,0.16083431)
\curveto(9.130576,0.095834315)(9.160576,-0.07916569)(9.165576,-0.18916568)
\curveto(9.170576,-0.2991657)(9.1455765,-0.4341657)(9.115577,-0.4591657)
\curveto(9.085576,-0.4841657)(9.0205765,-0.5841657)(8.985577,-0.6591657)
\curveto(8.950577,-0.73416567)(8.8955765,-0.8191657)(8.835576,-0.8491657)
}
\psline[linecolor=black, linewidth=0.04](8.595576,-1.1891657)(7.795576,-2.2891657)(7.8155766,-2.3091657)
\rput[bl](4.8755765,0.030834312){\LARGE{$=$}}
\end{pspicture}
}
\end{aligned}
\end{equation}

There is a braided monoidal category $\mathfrak{B}$ called the {\em braid groupoid}.
The objects are the natural numbers. Morphisms $m\xra{s} m$ are equivalence classes
of string diagrams involving only $m$ strings and nodes labelled by $c$ or $c^{-1}$;
these are called {\em braids on $m$ strings}.
The equation of \eqref{crossing/YB} gives two equivalent string diagrams for a morphism
$3\to 3$. The composition $\circ$ and tensor product $\ot$ are as for string diagrams;
in particular, on objects the tensor product is addition $m+n$ of natural numbers.
The braiding $m+n\xra{c_{m,n}} n+m$ takes $m+n$ vertical parallel strings and
places the first $m$ strings over the second $n$ strings.  

There is a symmetric monoidal category $\mathfrak{S}$ called the {\em permutation groupoid}.
It is as for $\mathfrak{B}$ except that whether the crossings are over or under is ignored so that morphisms $m\to m$ are permutations of $n$ symbols.  

If you take a braid $m\xra{s} m$ on $m$ strings and pull the bottom ends around, say to the right of the diagram, and splice them in order onto the top ends of the strings, we obtain what is called a {\em link}. This process is called the {\em Markov closure} of the braid. The number of
strings involved can reduce in the process. If the result is one string, we have a {\em knot}.

We study links by projecting them onto the plane taking note of under and over crossings
and making sure in the projection no two crossovers are projected to the same point. 
However, it is quite difficult to tell whether two knot projections actually represent the same 
link. So we look for numbers calculated from the projections which give the same result for
the same original 3-dimensional link: numerical invariants. No single number will completely
distinguish different links so we look at invariant polynomials which amount to having
a sequence of coefficient numbers. The Alexander polynomial \cite{Ale} stood as all we had
in this direction for over half a century. This polynomial could be calculated from
homology, an example of a functor in algebraic topology, so this is already some connection
with category theory.

In 1985, Vaughan Jones announced \cite{Jon} a revolutionary polynomial invariant for oriented links. This, and his work that led to it, earned the New Zealander a Fields Medal. The
consequences of this discovery were enormous and are ongoing. In the next issue of
the same journal and in the light of Jones' discovery, a six person paper \cite{HOMFLY} appeared: this paper combined four independent submissions with different proofs of a more general two-variable polynomial invariant for oriented links; it is called the {\em HOMFLY polynomial} from surnames of the authors. In particular, it showed that category theorists 
Freyd and Yetter were actively working on knot theory; their proof invoked term rewriting ideas as in $\lambda$-calculus, computer science and algebra. 

Soon after and published as \cite{Yetter1988}, Yetter constructed a monoidal category $\mathfrak{T}$ whose morphisms were {\em tangles} which include both braids and links. The strings in a braid can always be taken as progressing downwards but links inescapably involve turning back up at times.  More explicitly, the objects of $\mathfrak{T}$ are words $- + + - - +$ in symbols $+$ and $-$; morphisms are tangles such as
 \[
\xy
(0,20)*+!D{-}="t1";(10,20)*+!D{-}="t2";(20,20)*+!D{-}="t3";
(30,20)*+!D{+}="t4";(40,20)*+!D{+}="t5"; (50,20)*+!D{+}="t6";(60,20)*+!D{-}="t7";
(0,-10)*+!U{-}="b1";(10,-10)*+!U{+}="b2";(20,-10)*+!U{-}="b3";
"t3"+(3,-8)="a";
"t2"+(8,-15)="c";
"t3";"a"*{\hole} **\crv{"t3"+(0,-3)&"a"+(-1,2)} \POS?(.5)="b";
"t4";"b"*{\hole} **\crv{"t4"+(0,-1)&"b"+(1,1)} \POS?(.7)*\dir{>};
"t2";"c"*{\hole} **\crv{"t2"+(0,-2)&"c"+(-3,2)} \POS?(.8)="d" \POS?(.6)="e";
"b"*{\hole};"d"*{\hole} **\crv{"b"-(2,2)&"d"+(3,1)} \POS?(.5)="f";
"e"*{\hole};"f"*{\hole} **\crv{"f"};
"f"*{\hole};"a" **\crv{"f"&"a"+(-2,0)};
"e"+(-8,-2)="g";
"e"*{\hole};"g" **\crv{"e"+(1,0)&"g"+(2,1)} \POS?(.5)*\dir{<};
"t1";"g"*{\hole} **\crv{"t1"+(0,-3)&"g"+(0,3)} \POS?(.5)*\dir{<};
"d"+(-4,-4)="h";
"h"*{\hole};"b1" **\crv{"h"+(-4,0)&"b1"+(1,2)} \POS?(.5)="j" \POS?(.25)="k" \POS?(.65)*\dir{<};
"h"*{\hole};"c" **\crv{"h"+(2,0)&"c"+(-2,0)};
"d"*{\hole};"h" **\crv{"d"-(3,1)&"h"+(0,2)};
"g";"g"+(-5,-3)="i" **\crv{"g"+(-2,-1)&"i"+(2,2)};
"i";"i"*{\hole} **\crv{"i"-(8,8)&"i"+(-8,8)};
"i"*{\hole};"j"*{\hole} **\crv{"i"+(2,-2)&"j"+(-2,2)};
"g"*{\hole};"k"*{\hole} **\crv{"g"-(0,2)&"k"+(-3,1)};
"k"*{\hole};"b3" **\crv{"k"+(3,-1)&"b3"+(0,3)} \POS?(.3)="m" \POS?(.6)="l";
"j"*{\hole};"l"*{\hole} **\crv{"j"+(2,-2)&"l"+(-3,0)} \POS?(.4)="n";
"h";"m"*{\hole} **\crv{"h"-(0,2)&"m"+(2,2)};
"m"*{\hole};"n"*{\hole} **\crv{"m"-(2,2)&"n"+(1,2)};
"n"*{\hole};"b2" **\crv{"n"-(1,2)&"b2"+(0,2)};
"l"+(12,-1)="o";
"l"*{\hole};"o"*{\hole} **\crv{"l"+(2,0)&"o"+(-2,0)};
"c"*{\hole};"o" **\crv{"c"-(-3,2)&"o"+(-1,3)};
"o";"t6" **\crv{"o"+(4,-12)&"t6"-(0,5)} \POS?(.6)*\dir{<};
"a"+(6,-4)="p";
"a";"p"*{\hole} **\crv{"a"+(3,0)&"p"+(-2,2)};
"p"*{\hole};"o"*{\hole} **\crv{"p"+(2,-2)&"o"+(4,0)} \POS?(.5)="q";
"a"*{\hole};"q"*{\hole} **\crv{"a"+(1,-2)&"q"+(-4,0)} \POS?(.45)="r";
"q"*{\hole};"p" **\crv{"q"+(6,0)&"p"+(4,2)};
"p";"r"*{\hole} **\crv{"p"-(2,1)&"r"+(2,1)};
"r"*{\hole};"c" **\crv{"r"-(2,1)&"c"+(2,0)};
"t5";"t7" **i\crv{"t5"+(0,-10)&"t7"+(0,-10)} \POS?(.42)="s";
"t5";"s"*{\hole} **\crv{"t5"+(0,-5)&"s"-(5,0)};
"s"*{\hole};"t7" **\crv{"s"+(3,0)&"t7"+(0,-10)} \POS?(.5)*\dir{>}
\endxy
\]
Together with Freyd, in \cite{FreydYetter1989}, they recognized not only the presence of a braiding for $\mathfrak{T}$ but also that the backtracking of the strings was explained by
the presence of duals for each object; see Example~\ref{monoidaldual}. To see this, look back at the ``snake'' equations \eqref{snakes} and observe that, if we replace the unit node labelled $d$ by a cap $\cap$ and the counit node labelled $e$ by a cup $\cup$, the equations
become geometrically obvious. I asked my student Mei Chee Shum to look into a freeness
property for the category of tangles along the lines of a theorem proved by Kelly-Laplaza
for compact monoidal categories. Her thesis \cite{Shu} proved a universal property of the 
{\em tortile} monoidal category $\mathfrak{T^{\sim}}$ of tangles on ribbons rather than strings;
the objects are as for $\mathfrak{T}$.
Tortility \cite{BTC} includes a braiding, a twist operation, and existence of duals, subject to axioms. 
The universal property means that any object $A$ of a tortile monoidal category $\CA$ determines a
functor $\hat{A} : \mathfrak{T^{\sim}}\to \CA$ which preserves the tensor product, braiding and twist,
and whose value at $+$ is the object $A$. 
  
Meanwhile, mathematical physicists and low-dimensional topologists (see \cite{Dri1985, FRT, Turaev1988, Seg, Wit}) were studying the quantum inverse scattering method, Yang-Baxter equations, 
quantum groups, conformal field theory, and topological quantum field theory. 
The Introduction of \cite{40} has further detail on Tannaka duality and quantum groups.

A Yang-Baxter operator on an object $A$ of a monoidal category $\CA$ 
is an invertible morphism $y : A\ot A \to A\ot A$ such that the equation of \eqref{crossing/YB} holds
where the three strings are all labelled $A$ and the crossings depict $y$. 
If $\CA$ is braided, each $c_{A,A} : A\ot A\to A\ot A$ is a Yang-Baxter operator on $A$.
A Yang-Baxter operator $y$ on $A$ is {\em tortile} when it is equipped with a twist isomorphism $z : A \to A$
and $A$ has a dual subject to some compatibility conditions.
 In \cite{TYBO}, we prove that $\mathfrak{T^{\sim}}$ has another universal property: each tortile Yang-Baxter operator $(y,z)$ in any monoidal category $\CA$ determines a unique (up to equivalence) functor $\mathfrak{T^{\sim}}\to A$ which preserves (up to isomorphism) the tensor product and takes $c_{+,+} : ++\to ++$ to $y$ and the ribbon twist to $z$. This property fascilitates the construction of representations of the ribbon tangle category $\mathfrak{T^{\sim}}$.

\section{Frobenius monoids}\label{Frobenius}

One of the great powers of monoidal category theory is that a structure that 
can be expressed in terms of string diagrams in one monoidal category can be sought in others.
If we are really lucky, the free monoidal category containing that structure might have an
explicit description. We already have examples of that: $\mathfrak{B}$ is the free monoidal
category containing an object bearing a Yang-Baxter, and $\mathfrak{T}^{\sim}$ is the free monoidal
category containing an object bearing a tortile Yang-Baxter. 

A more basic example is that of {\em monoid in a monoidal category $\CV$}. A monoid is an object $A$
equipped with morphisms $A\ot A\xra{m} A$ and $I\xra{i} A$ satisfying the associativity and identity
axioms depicted in diagram~\eqref{monoidaxx} where all strings are labelled $A$.

\begin{equation}\label{monoidaxx}
\begin{aligned}
\psscalebox{0.5 0.6} % Change this value to rescale the drawing.
{
\begin{pspicture}(0,-2.7442338)(23.45826,2.7442338)
\pscircle[linecolor=black, linewidth=0.04, dimen=outer](20.36826,1.0460731){0.3}
\pscircle[linecolor=black, linewidth=0.04, dimen=outer](11.368259,1.0060731){0.3}
\pscircle[linecolor=black, linewidth=0.04, dimen=outer](7.288259,1.0060731){0.3}
\pscircle[linecolor=black, linewidth=0.04, dimen=outer](5.708259,-1.1939269){0.3}
\pscircle[linecolor=black, linewidth=0.04, dimen=outer](2.0082593,-1.1139269){0.3}
\pscircle[linecolor=black, linewidth=0.04, dimen=outer](0.92825925,1.0260731){0.3}
\pscircle[linecolor=black, linewidth=0.04, dimen=outer](12.928259,-1.133927){0.3}
\pscircle[linecolor=black, linewidth=0.04, dimen=outer](18.64826,-1.1139269){0.3}
\rput[bl](18.51826,-1.2239269){\footnotesize{$m$}}
\rput[bl](12.798259,-1.2239269){\footnotesize{$m$}}
\rput[bl](7.1182594,0.89607304){\footnotesize{$m$}}
\rput[bl](5.5982594,-1.283927){\footnotesize{$m$}}
\rput[bl](0.81825924,0.93607306){\footnotesize{$m$}}
\rput[bl](1.8982593,-1.2239269){\footnotesize{$m$}}
\rput[bl](20.31826,0.95607305){\footnotesize{$i$}}
\rput[bl](11.278259,0.9160731){\footnotesize{$i$}}
\psline[linecolor=black, linewidth=0.04](0.018259238,2.736073)(0.69825923,1.216073)
\psline[linecolor=black, linewidth=0.04](2.0382593,2.716073)(1.1182592,1.2560731)
\psline[linecolor=black, linewidth=0.04](0.9582592,0.7360731)(1.7982593,-0.8839269)
\psline[linecolor=black, linewidth=0.04](2.1582592,-0.8839269)(3.3982592,2.736073)
\psline[linecolor=black, linewidth=0.04](1.9782592,-1.403927)(1.9982592,-2.743927)
\psline[linecolor=black, linewidth=0.04](4.918259,2.716073)(5.5782595,-0.9639269)
\psline[linecolor=black, linewidth=0.04](7.2782593,0.71607304)(5.898259,-0.98392695)
\psline[linecolor=black, linewidth=0.04](6.2782593,2.696073)(7.1182594,1.196073)
\psline[linecolor=black, linewidth=0.04](7.478259,1.196073)(8.398259,2.696073)
\psline[linecolor=black, linewidth=0.04](5.6782594,-1.4839269)(5.6982594,-2.743927)
\psline[linecolor=black, linewidth=0.04](11.37826,0.7360731)(12.758259,-0.92392695)
\psline[linecolor=black, linewidth=0.04](13.078259,-0.86392695)(13.078259,2.656073)
\psline[linecolor=black, linewidth=0.04](12.938259,-1.423927)(12.938259,-2.723927)
\psline[linecolor=black, linewidth=0.04](18.418259,2.556073)(18.45826,-0.86392695)
\psline[linecolor=black, linewidth=0.04](18.758259,-0.86392695)(20.178259,0.8560731)
\psline[linecolor=black, linewidth=0.04](18.598259,-1.3639269)(18.578259,-2.743927)
\psline[linecolor=black, linewidth=0.04](15.758259,2.636073)(15.758259,-2.743927)
\rput[bl](17.11826,-0.6439269){\LARGE{$=$}}
\rput[bl](14.198259,-0.68392694){\LARGE{$=$}}
\rput[bl](3.7582593,-0.66392696){\LARGE{$=$}}
\rput[bl](9.298259,-0.78392696){\LARGE{$,$}}
\end{pspicture}
}
\end{aligned}
\end{equation}
The free monoidal category containing a monoid is denote by $\Delta$ and is called the 
{\em augmented simplex category}. The objects are the ordinals $\mathbf{n} = \{0, 1, 2, \dots , n-1\}$, each
$\mathbf{n}$ abstracting $n$ points in general position in $\mathbb{R}^{n-1}$ forming an $n$-simplex
(a 3-simplex is a triangle and a 4-simplex is a tetrahedron). In this case the morphisms are functions: 
they are the monotone non-decreasing functions
$\mathbf{n}\xra{\xi}\mathbf{m}$ between ordinals. On objects, the tensor product is addition: 
$\mathbf{n}\ot \mathbf{n'} = \mathbf{n}+ \mathbf{n'}$. We have the monoid $\mathbf{1}$ with
multiplication the unique function $\mathbf{2}\to \mathbf{1}$ and identity the unique function
 $\mathbf{0}\to \mathbf{1}$. Every morphism can be obtained from these 
 operations using composition and tensoring. The simplicial sets mentioned at the end of Section~\ref{Qt}
 are functors from $\Delta^{\mathrm{op}}$ to $\mathrm{Set}$ which take $\mathbf{0}$ to a
 one-element set.  
 
 The concept of Frobenius algebra is classical. It is a good example of a structure in the category of vector spaces that can be expessed using the monoidal category structure. A monoid $A$ in a monoidal category
 $\CV$ is {\em Frobenius} when it is equipped with a morphism $A\ot A\xra{r}I$ which is a counit for
 an adjunction $A\dashv A$ and satisfies the condition \eqref{Frobax} where all strings are labelled $A$.
 \begin{equation}\label{Frobax}
\begin{aligned}
\psscalebox{0.5 0.6} % Change this value to rescale the drawing.
{
\begin{pspicture}(0,-2.2646327)(11.956815,2.2646327)
\pscircle[linecolor=black, linewidth=0.04, dimen=outer](2.5968146,-1.8446326){0.4}
\pscircle[linecolor=black, linewidth=0.04, dimen=outer](6.8568144,-1.8646326){0.4}
\pscircle[linecolor=black, linewidth=0.04, dimen=outer](8.676815,0.1753674){0.4}
\pscircle[linecolor=black, linewidth=0.04, dimen=outer](1.3568146,0.1953674){0.4}
\rput[bl](8.576815,0.07536739){\footnotesize{$m$}}
\rput[bl](1.2368146,0.07536739){\footnotesize{$m$}}
\rput[bl](6.7368145,-1.9646326){\footnotesize{$r$}}
\rput[bl](2.4568145,-1.9246325){\footnotesize{$r$}}
\rput[bl](4.8368144,-0.14463261){\LARGE{$=$}}
\psline[linecolor=black, linewidth=0.04](1.1368146,0.4953674)(0.016814575,2.2353673)
\psline[linecolor=black, linewidth=0.04](1.6568146,0.4753674)(2.3968146,2.2353673)
\psline[linecolor=black, linewidth=0.04](1.4168146,-0.1846326)(2.3568146,-1.5646327)
\psline[linecolor=black, linewidth=0.04](2.7368145,-1.5046326)(4.1168146,2.2553673)
\psline[linecolor=black, linewidth=0.04](6.5968146,-1.5446326)(5.7568145,2.2553673)
\psline[linecolor=black, linewidth=0.04](7.0568147,-1.5246326)(8.676815,-0.20463261)
\psline[linecolor=black, linewidth=0.04](8.416815,0.4553674)(7.3568144,2.2353673)
\psline[linecolor=black, linewidth=0.04](8.936814,0.4553674)(9.876815,2.2553673)
\end{pspicture}
}
\end{aligned}
\end{equation}
It turns out that there is a lovely geometric model $\mathrm{Cob_2}$ of the 
free monoidal category containing a Frobenius monoid. The objects are the natural numbers.
A morphism $n\xra{\Gamma} m$ is a surface in 3D whose boundary consists of $m+n$ circles with $n$ of those circles
in a horizontal plane as domain and $m$ in a parallel horizontal plane below as codomain. For example, Figure 1 depicts a morphism $3\to 2$. 
\begin{figure} [htbp]
\centering
% \includegraphics [width=0.4\textwidth]{CobordEx1.eps}
\includegraphics[scale=0.2, natwidth=480, natheight=765]{CobordEx1.eps}
\vspace{-5em} % Adjust -1em to make the space smaller or larger
\caption {Cobordism}
% \label{fig:my_pdf}
\end{figure}
The Frobenius structure $m, i, r$ on the object $1$ are shown in
Figure 2 with the axioms in Figure 3. Composition is by vertical joining and 
tensor product is by horizontal placement. The category of functors $\mathrm{Cob_2}\xra{F} \CV$ into
a monoidal category $\CV$, which coherently preserve tensor product up to isomorphism, is equivalent to
the category of Frobenius monoids in $\CV$. Such a functor $F$ amounts to a 2-dimensional Topological Quantum Field Theory; see \cite{Koc}.

Higher dimensional Topological Quantum Field Theories also have categorical constructions (see \cite{TuraevViro1992, Lyu, TurVirel}). They have implications, not only for theoretical physics, but
for the mathematics of invariants for surfaces and higher-dimensional spaces.

% Side by side figures 
\begin{figure} [htbp]
\begin{minipage}[c]{0.4\linewidth}
\includegraphics[width=\linewidth]{CobordStruct1.eps}
\caption{Structure}
\end{minipage}
\hfill
\begin{minipage}[c]{0.4\linewidth}
\includegraphics[width=\linewidth]{CobordAxioms1.eps}
\caption{Axioms}
\end{minipage}
\end{figure}

\begin{center}
--------------------------------------------------------
\end{center}

\newpage
\appendix
\begin{thebibliography}{000}

\bibitem{Ale} J.W. Alexander, 
\textit{Topological invariants of knots and links}, 
Transactions of the American Math. Soc. \textbf{30} (1928) 275--306.

\bibitem{BaeDol1} J. Baez and J. Dolan, 
\textit{Higher dimensional algebra and topological quantum field theory}, 
Journal of Mathematical Physics \textbf{36} (1995) 6073--6105.

\bibitem{BaeDol2} J. Baez and J. Dolan, 
\textit{From finite sets to Feynman diagrams}, 
in: ``Mathematics unlimited -- 2001 and beyond''
(Springer-Verlag, Berlin, 2001) 29--50.

\bibitem{Ben} J. B\'enabou, \textit{Introduction to bicategories}, 
Lecture Notes in Math. \textbf{47} (Springer-Verlag, 1967) 1--77.

\bibitem{BluCocSee} R.F. Blute, J.R.B. Cockett and R.A.G. Seely, 
\textit{Differential categories}, 
Mathematical Structures in Computer Science \textbf{16(06)} (2006) 1049--1083.

\bibitem{BluCocSee2} R.F. Blute, J.R.B. Cockett and R.A.G. Seely, 
\textit{Cartesian differential categories}, 
Theory and Applications of Categories \textbf{23} (2009) 622--672.

\bibitem{BluCocLemSee}R. Blute, J.R.B. Cockett, J.-S.P. Lemay and R.A.G. Seely, \textit{Differential categories revisited}, 
Applied Categorical Structures  \textbf{28} (2020) 171--235. 

\bibitem{Che} C. Chevalley, 
\textit{Theory of Lie Groups}, (Princeton University Press, 1946).

\bibitem{CocGal} J.R.B. Cockett and J. Gallagher, 
\textit{Categorical models of the differential $\lambda$-calculus}, 
Mathematical Structures in Computer Science \textbf{29(10)} (2019) 1513--1555.

\bibitem{CoeGog} B. Coecke and S. Gogioso, 
\textit{Quantum in Pictures}, 
(Quantinuum, 2022; ISBN 978-1-7392147-1-5).

\bibitem{Del} P. Deligne, \textit{Cat\'egories tannakiennes}, 
in: ``Grothendieck's Festschrift'', Progress in Math. (Birkh\"auser, 1990).

\bibitem{DelMil} P. Deligne and J.S. Milne, 
 \textit{Tannakian categories, Hodge Cocycles Motives and Shimura Varieties},
Lecture Notes in Math. \textbf{900} (Springer-Verlag, Berlin, 1982) 101--228.

\bibitem{DopRob} S. Doplicher and J.E. Roberts, 
\textit{A new duality theory for compact groups}, Inventiones Math. \textbf{98} (1989)
157--218.

\bibitem{Dri1985} V.G. Drinfel'd, 
\textit{Hopf algebras and the quantum Yang-Baxter equation},
Dokl. Akad. Nauk SSSR \textbf{283(5)} (1985) 1060--1064. (Russian)

\bibitem{Dri} V.G. Drinfel'd, 
\textit{Quantum groups}, 
Proceedings of the International Congress of Mathematicians at
Berkeley, California, U.S.A. 1986 (1987) 798--820.

\bibitem{Ehres} C. Ehresmann, 
\textit{Cat\'egories doubles et cat\'egories structur\'ees}, 
Comptes Rendus de l'Acad\'emie des Sciences Paris \textbf{256} (1963) 1198--1201.

\bibitem{EhrReg} T. Ehrhard and L. Regnier,
\textit{The differential lambda-calculus},
Theoretical Computer Science \textbf{309} (2003) 1--41.

\bibitem{EilKel1966} S. Eilenberg and G.M. Kelly, 
\textit{Closed categories}, 
in: ``Proceedings of the Conference on Categorical Algebra (La Jolla, 1965)'', (Springer-Verlag,1966) 421--562.

\bibitem{EilMac} S. Eilenberg and S. Mac Lane,
\textit{General Theory of Natural Equivalences},
Transactions of the American Math. Soc. \textbf{58(2)} (1945) 231--294; \url{https://www.jstor.org/stable/1990284} 
\label{EilMac}

\bibitem{FRT} L.D. Faddeev, N.Yu. Reshetikhin and L.A. Takhtajan, 
\textit{Quantization of Lie algebras and Lie groups},
LOMI Preprints (Leningrad 1987); Algebra i Analiz \textbf{1:1} (1989, in Russian).

\bibitem{HOMFLY} P. Freyd, D.N. Yetter; W.B.R. Lickorish, K.C. Millet; J. Hoste; and A. Ocneanu, 
\textit{A new polynomial invariant of knots and links}, 
Bulletin of the American Math. Soc. \textbf{12} (1985) 239--246.

\bibitem{FreydYetter1989} P. Freyd and D. Yetter, 
\textit{Braided compact closed categories with applications to low dimensional topology}, 
Advances in Math. \textbf{77} (1989) 156--182.

\bibitem{Gro} A. Grothendieck, 
\textit{Esquisse d'un programme}, 
Research proposal to join the CNRS (1984);
\newline \url{fr.wikipedia.org/wiki/Esquisse_d\% 27un_programme}

\bibitem{Jon} V. Jones, 
\textit{A polynomial invariant for knots via von Neumann algebras}, 
Bulletin American Math. Soc. \textbf{12} (1985) 103--111.

\bibitem{ix} A. Joyal and R. Street, 
\textit{Braided monoidal categories}, 
Macquarie Mathematics Reports 850067 (1985).

\bibitem{TYBO} A. Joyal and R. Street, 
\textit{Tortile Yang-Baxter operators in tensor categories}, 
Journal of Pure and Applied Algebra \textbf{71} (1991) 43--51. 

\bibitem{GTC} A. Joyal and R. Street, 
\textit{The geometry of tensor calculus I}, 
Advances in Math. \textbf{88} (1991) 55--112.

\bibitem{40} A. Joyal and R. Street, 
\textit{An introduction to Tannaka duality and quantum groups}, Lecture Notes in Math. \textbf{1488} (Springer-Verlag Berlin, Heidelberg 1991) 411--492.

\bibitem{BTC}  A. Joyal and R. Street, 
\textit{Braided tensor categories}, 
Advances in Math. \textbf{102} (1993) 20--78.

\bibitem{Kan1958} D.M. Kan, \textit{Adjoint functors}, 
Transactions of the American Math. Soc. \textbf{87} (1958) 294--329. \label{Kan1958}

\bibitem{Koc} J. Kock,
\textit{Frobenius Algebras and 2D Topological Quantum Field Theories},
London Math. Society Student Texts \textbf{59} (Cambridge University Press, 2003).

\bibitem{KelBk} G. Max Kelly, 
\textit{Basic concepts of enriched category theory}, 
London Math. Soc. Lecture Note Series \textbf{64} (Cambridge University Press, Cambridge, 1982); also, Reprints in Theory and Applications of Categories, \textbf{10} (2005) 1--136. 

\bibitem{7} G.M. Kelly and R. Street, 
\textit{Review of the elements of 2-categories}, 
Lecture Notes in Math. \textbf{420} (1974) 75--103.

\bibitem{Kit} A.Yu. Kitaev
\textit{Fault-tolerant quantum computation by anyons},
Annals of Physics \textbf{303} (2003) 2--30.

\bibitem{Knu} D. Knuth,
\textit{The Art of Computer Programming, Volumes 1-4A Boxed Set}, 
(Third Edition, Addison-Wesley, 2011) 3168pp.

\bibitem{Lam} J. Lambek, 
\textit{Deductive systems and categories I}, 
Journal of Mathematical Systems Theory \textbf{2} (1968) 278--318.
\label{Lam}

\bibitem{LamSco} J. Lambek and P.J. Scott, 
\textit{Introduction to higher order categorical logic}, 
Cambridge studies in advanced mathematics \textbf{7} (CUP, 1986).
\label{LamSco}

\bibitem{LiWil} A. Li and D.J. Williamson,
\textit{A classification of translation-invariant quantum codes in any dimension},
arXiv:26081v1 [quant-ph] (21 Aug 2026) 7pp.

\bibitem{Lyu} V.V. Lyubashenko, 
\textit{Invariants of 3-manifolds and projective representations of mapping class groups 
via quantum groups at roots of unity} 
Communications in Mathematical Physics 
\textbf{172} (1995) 467--516.

\bibitem{ML1963} S. Mac Lane, 
\textit{Natural associativity and commutativity}, 
Rice University Studies \textbf{49} (1963) 28--46.

\bibitem{CWM} S. Mac Lane, 
\textit{Categories for the Working Mathematician}, 
Graduate Texts in Mathematics \textbf{5} (Springer-Verlag, 1971).

\bibitem{Man} C.R. Mann,
\textit{The connection between equivalence of proofs and cartesian closed categories},
Proc. of the London Math. Soc. \textbf{31} (1975) 289--310.

\bibitem{Pen1971} R. Penrose, 
\textit{Applications of negative dimensional tensors}, 
in: ``Combinatorial Mathematics and its Applications'' (Academic Press 1971) 221--244.

\bibitem{Pow} A.J. Power, 
\textit{An n-categorical pasting theorem}, 
Lecture Notes in Math. \textbf{1488} (Springer-Verlag, 1991) 326--358.

\bibitem{Pra} D. Prawitz,
\textit{Ideas and results in proof theory},
In Fenstad (1971) 235--307.

\bibitem{ResTur} N.Yu. Reshetikhin and V.G. Turaev, 
\textit{Ribbon graphs and their invariants derived from quantum groups}, 
Comm. Math. Phys. \textbf{127(1)} (1990) 1--26.

\bibitem{Rob} J.E. Roberts, 
\textit{Mathematical aspects of local cohomology}, 
in: Proceedings of the Colloquium on Operator Algebras 
and their Application to Mathematical Physics (Marseille, 1977).

\bibitem{Rob} J.E. Roberts, 
\textit{Net Cohomology and its Applications to Field Theory},
Presented at: Bielefeld Encounters in Physics and Mathematics II: 
``Quantum Fields -- Algebras, Processes'' (December 1--9, 1978).

\bibitem{Ros} J. Rosick\'y, 
\textit{Abstract tangent functors}, 
Diagrammes \textbf{12(3)} (1984).

\bibitem{Saa} N. Saavedra Rivano, \textit{Cat\'egories Tannakiennes}, 
Lecture Notes in Math. \textbf{265} (Springer-Verlag, Berlin 1972).

\bibitem{See} R.A.G. Seely,
\textit{Hyperdoctrines and natural deduction},
(PhD Thesis, University of Cambridge, U.K., 1977).

\bibitem{Seg} G. Segal,
\textit{Two-dimensional conformal field theories and modular functors}, 
in: ``IXth International Congress on Mathematical Physics (Swansea, 1988)'', 
(Hilger, Bristol, 1989) 22--37.

\bibitem{Shu} M.C. Shum,
\textit{Tortile Tensor Categories} (PhD thesis, Macquarie University, November 1989); 
Journal of Pure and Applied Algebra \textbf{93} (1994) 57--110.

\bibitem{3} R. Street,
\textit{The formal theory of monads}, 
Journal of Pure and Applied Algebra \textbf{2} (1972) 149--168.

\bibitem{30} R. Street,
\textit{The algebra of oriented simplexes}, 
Journal of Pure and Applied Algebra \textbf{49} (1987) 283--335.

\bibitem{113} R. Street,
\textit{Monoidal categories in, and linking, geometry and algebra}, 
Bulletin of the Belgian Mathematical Society - Simon Stevin \textbf{19(5)} (2012) 769--821.

\bibitem{Tan} T. Tannaka,
\textit{ \"Uber den Dualit\"atssatz der nichtkommutativen topologischen Gruppen}, 
T\^ohoku Math. Journal \textbf{45} (1939) 1--12.

\bibitem{Turaev1988} V.G. Turaev, \textit{The Yang-Baxter equation and invariants of links}, Inventiones Mathematicae \textbf{92} (1988) 527--553.\label{Turaev1988}

\bibitem{TurVirel} V.G. Turaev and A. Virelizier, \textit{On two approaches to 3-dimensional TQFTs}, arXiv:1006.3501 (June 2010).\label{TurVirel}

\bibitem{TuraevViro1992} V.G. Turaev and O. Viro, \textit{State sum invariants of 3-manifolds and quantum $6j$-symbols}, Topology \textbf{31} (1992) 865--902.\label{TuraevViro1992}

\bibitem{VerO} D. Verity,
\textit{``Complicial sets characterising the simplicial nerves of strict $\omega$-categories''},
Memoirs of the American Math. Soc. \textbf{905} (American Mathematical Society, 2008)

\bibitem{Ver} D. Verity,
\textit{Weak complicial sets I, basic homotopy theory},
Advances in Math. \textbf{219} (2008) 1081--1149.

\bibitem{Wey1918} H. Weyl, 
\textit{``Raum$\ \cdot \ $Zeit$\ \cdot \ $Materie: Vorlesungen \"uber allgemeine Relativit\"atstheorie''}
(Springer, Berlin, 1918) 234pp.

\bibitem{Wey1927} H. Weyl, 
\textit{Quantenmechanik und Gruppentheorie}, 
Zeitschrift f\"ur Physik \textbf{46(1-2)} (1927) 1--46.

\bibitem{Wit} E. Witten, \textit{Topological quantum field theory}, 
Communications in Mathematical Physics \textbf{117} (1988) 353--386.\label{Wit}

\bibitem{Yetter1988} D.N. Yetter, \textit{Markov algebras}, 
Contemporary Math. \textbf{78} (American Math. Soc., 1988) 705--730.\label{Yetter1988}

\end{thebibliography}

\end{document}